\documentclass{report}

%%%%%%%%%%%%%%%%%%%%%%%%%%%%%%%%%
% PACKAGE IMPORTS
%%%%%%%%%%%%%%%%%%%%%%%%%%%%%%%%%


\usepackage[tmargin=2cm,rmargin=1in,lmargin=1in,margin=0.85in,bmargin=2cm,footskip=.2in]{geometry}
\usepackage{amsmath,amsfonts,amsthm,amssymb,mathtools}
\usepackage[varbb]{newpxmath}
\usepackage{xfrac}
\usepackage[makeroom]{cancel}
\usepackage{mathtools}
\usepackage{bookmark}
\usepackage{enumitem}
\usepackage{hyperref,theoremref}
\hypersetup{
	pdftitle={Assignment},
	colorlinks=true, linkcolor=doc!90,
	bookmarksnumbered=true,
	bookmarksopen=true
}
\usepackage[most,many,breakable]{tcolorbox}
\usepackage{xcolor}
\usepackage{varwidth}
\usepackage{varwidth}
\usepackage{etoolbox}
%\usepackage{authblk}
\usepackage{nameref}
\usepackage{multicol,array}
\usepackage{tikz-cd}
\usepackage[ruled,vlined,linesnumbered]{algorithm2e}
\usepackage{comment} % enables the use of multi-line comments (\ifx \fi) 
\usepackage{import}
\usepackage{xifthen}
\usepackage{pdfpages}
\usepackage{transparent}
\usepackage[utf8]{inputenc}
\usepackage{fontawesome}


\newcommand\mycommfont[1]{\footnotesize\ttfamily\textcolor{blue}{#1}}
\SetCommentSty{mycommfont}
\newcommand{\incfig}[1]{%
    \def\svgwidth{\columnwidth}
    \import{./figures/}{#1.pdf_tex}
}

\usepackage{tikzsymbols}
\renewcommand{\qedsymbol}{\faWhatsapp}


%\usepackage{import}
%\usepackage{xifthen}
%\usepackage{pdfpages}
%\usepackage{transparent}


%%%%%%%%%%%%%%%%%%%%%%%%%%%%%%
% SELF MADE COLORS
%%%%%%%%%%%%%%%%%%%%%%%%%%%%%%



\definecolor{myg}{RGB}{56, 140, 70}
\definecolor{myb}{RGB}{45, 111, 177}
\definecolor{myr}{RGB}{199, 68, 64}
\definecolor{mytheorembg}{HTML}{F2F2F9}
\definecolor{mytheoremfr}{HTML}{00007B}
\definecolor{mylenmabg}{HTML}{FFFAF8}
\definecolor{mylenmafr}{HTML}{983b0f}
\definecolor{mypropbg}{HTML}{f2fbfc}
\definecolor{mypropfr}{HTML}{191971}
\definecolor{myexamplebg}{HTML}{F2FBF8}
\definecolor{myexamplefr}{HTML}{88D6D1}
\definecolor{myexampleti}{HTML}{2A7F7F}
\definecolor{mydefinitbg}{HTML}{E5E5FF}
\definecolor{mydefinitfr}{HTML}{3F3FA3}
\definecolor{notesgreen}{RGB}{0,162,0}
\definecolor{myp}{RGB}{197, 92, 212}
\definecolor{mygr}{HTML}{2C3338}
\definecolor{myred}{RGB}{127,0,0}
\definecolor{myyellow}{RGB}{169,121,69}
\definecolor{myexercisebg}{HTML}{F2FBF8}
\definecolor{myexercisefg}{HTML}{88D6D1}


%%%%%%%%%%%%%%%%%%%%%%%%%%%%
% TCOLORBOX SETUPS
%%%%%%%%%%%%%%%%%%%%%%%%%%%%

\setlength{\parindent}{1cm}
%================================
% THEOREM BOX
%================================

\tcbuselibrary{theorems,skins,hooks}
\newtcbtheorem[number within=section]{Theorem}{Theorem}
{%
	enhanced,
	breakable,
	colback = mytheorembg,
	frame hidden,
	boxrule = 0sp,
	borderline west = {2pt}{0pt}{mytheoremfr},
	sharp corners,
	detach title,
	before upper = \tcbtitle\par\smallskip,
	coltitle = mytheoremfr,
	fonttitle = \bfseries\sffamily,
	description font = \mdseries,
	separator sign none,
	segmentation style={solid, mytheoremfr},
}
{th}

\tcbuselibrary{theorems,skins,hooks}
\newtcbtheorem[number within=chapter]{theorem}{Theorem}
{%
	enhanced,
	breakable,
	colback = mytheorembg,
	frame hidden,
	boxrule = 0sp,
	borderline west = {2pt}{0pt}{mytheoremfr},
	sharp corners,
	detach title,
	before upper = \tcbtitle\par\smallskip,
	coltitle = mytheoremfr,
	fonttitle = \bfseries\sffamily,
	description font = \mdseries,
	separator sign none,
	segmentation style={solid, mytheoremfr},
}
{th}


\tcbuselibrary{theorems,skins,hooks}
\newtcolorbox{Theoremcon}
{%
	enhanced
	,breakable
	,colback = mytheorembg
	,frame hidden
	,boxrule = 0sp
	,borderline west = {2pt}{0pt}{mytheoremfr}
	,sharp corners
	,description font = \mdseries
	,separator sign none
}

%================================
% Corollery
%================================
\tcbuselibrary{theorems,skins,hooks}
\newtcbtheorem[number within=section]{Corollary}{Corollary}
{%
	enhanced
	,breakable
	,colback = myp!10
	,frame hidden
	,boxrule = 0sp
	,borderline west = {2pt}{0pt}{myp!85!black}
	,sharp corners
	,detach title
	,before upper = \tcbtitle\par\smallskip
	,coltitle = myp!85!black
	,fonttitle = \bfseries\sffamily
	,description font = \mdseries
	,separator sign none
	,segmentation style={solid, myp!85!black}
}
{th}
\tcbuselibrary{theorems,skins,hooks}
\newtcbtheorem[number within=chapter]{corollary}{Corollary}
{%
	enhanced
	,breakable
	,colback = myp!10
	,frame hidden
	,boxrule = 0sp
	,borderline west = {2pt}{0pt}{myp!85!black}
	,sharp corners
	,detach title
	,before upper = \tcbtitle\par\smallskip
	,coltitle = myp!85!black
	,fonttitle = \bfseries\sffamily
	,description font = \mdseries
	,separator sign none
	,segmentation style={solid, myp!85!black}
}
{th}


%================================
% LENMA
%================================

\tcbuselibrary{theorems,skins,hooks}
\newtcbtheorem[number within=section]{Lenma}{Lenma}
{%
	enhanced,
	breakable,
	colback = mylenmabg,
	frame hidden,
	boxrule = 0sp,
	borderline west = {2pt}{0pt}{mylenmafr},
	sharp corners,
	detach title,
	before upper = \tcbtitle\par\smallskip,
	coltitle = mylenmafr,
	fonttitle = \bfseries\sffamily,
	description font = \mdseries,
	separator sign none,
	segmentation style={solid, mylenmafr},
}
{th}

\tcbuselibrary{theorems,skins,hooks}
\newtcbtheorem[number within=chapter]{lenma}{Lenma}
{%
	enhanced,
	breakable,
	colback = mylenmabg,
	frame hidden,
	boxrule = 0sp,
	borderline west = {2pt}{0pt}{mylenmafr},
	sharp corners,
	detach title,
	before upper = \tcbtitle\par\smallskip,
	coltitle = mylenmafr,
	fonttitle = \bfseries\sffamily,
	description font = \mdseries,
	separator sign none,
	segmentation style={solid, mylenmafr},
}
{th}


%================================
% PROPOSITION
%================================

\tcbuselibrary{theorems,skins,hooks}
\newtcbtheorem[number within=section]{Prop}{Proposition}
{%
	enhanced,
	breakable,
	colback = mypropbg,
	frame hidden,
	boxrule = 0sp,
	borderline west = {2pt}{0pt}{mypropfr},
	sharp corners,
	detach title,
	before upper = \tcbtitle\par\smallskip,
	coltitle = mypropfr,
	fonttitle = \bfseries\sffamily,
	description font = \mdseries,
	separator sign none,
	segmentation style={solid, mypropfr},
}
{th}

\tcbuselibrary{theorems,skins,hooks}
\newtcbtheorem[number within=chapter]{prop}{Proposition}
{%
	enhanced,
	breakable,
	colback = mypropbg,
	frame hidden,
	boxrule = 0sp,
	borderline west = {2pt}{0pt}{mypropfr},
	sharp corners,
	detach title,
	before upper = \tcbtitle\par\smallskip,
	coltitle = mypropfr,
	fonttitle = \bfseries\sffamily,
	description font = \mdseries,
	separator sign none,
	segmentation style={solid, mypropfr},
}
{th}


%================================
% CLAIM
%================================

\tcbuselibrary{theorems,skins,hooks}
\newtcbtheorem[number within=section]{claim}{Claim}
{%
	enhanced
	,breakable
	,colback = myg!10
	,frame hidden
	,boxrule = 0sp
	,borderline west = {2pt}{0pt}{myg}
	,sharp corners
	,detach title
	,before upper = \tcbtitle\par\smallskip
	,coltitle = myg!85!black
	,fonttitle = \bfseries\sffamily
	,description font = \mdseries
	,separator sign none
	,segmentation style={solid, myg!85!black}
}
{th}



%================================
% Exercise
%================================

\tcbuselibrary{theorems,skins,hooks}
\newtcbtheorem[number within=section]{Exercise}{Exercise}
{%
	enhanced,
	breakable,
	colback = myexercisebg,
	frame hidden,
	boxrule = 0sp,
	borderline west = {2pt}{0pt}{myexercisefg},
	sharp corners,
	detach title,
	before upper = \tcbtitle\par\smallskip,
	coltitle = myexercisefg,
	fonttitle = \bfseries\sffamily,
	description font = \mdseries,
	separator sign none,
	segmentation style={solid, myexercisefg},
}
{th}

\tcbuselibrary{theorems,skins,hooks}
\newtcbtheorem[number within=chapter]{exercise}{Exercise}
{%
	enhanced,
	breakable,
	colback = myexercisebg,
	frame hidden,
	boxrule = 0sp,
	borderline west = {2pt}{0pt}{myexercisefg},
	sharp corners,
	detach title,
	before upper = \tcbtitle\par\smallskip,
	coltitle = myexercisefg,
	fonttitle = \bfseries\sffamily,
	description font = \mdseries,
	separator sign none,
	segmentation style={solid, myexercisefg},
}
{th}

%================================
% EXAMPLE BOX
%================================

\newtcbtheorem[number within=section]{Example}{Example}
{%
	colback = myexamplebg
	,breakable
	,colframe = myexamplefr
	,coltitle = myexampleti
	,boxrule = 1pt
	,sharp corners
	,detach title
	,before upper=\tcbtitle\par\smallskip
	,fonttitle = \bfseries
	,description font = \mdseries
	,separator sign none
	,description delimiters parenthesis
}
{ex}

\newtcbtheorem[number within=chapter]{example}{Example}
{%
	colback = myexamplebg
	,breakable
	,colframe = myexamplefr
	,coltitle = myexampleti
	,boxrule = 1pt
	,sharp corners
	,detach title
	,before upper=\tcbtitle\par\smallskip
	,fonttitle = \bfseries
	,description font = \mdseries
	,separator sign none
	,description delimiters parenthesis
}
{ex}

%================================
% DEFINITION BOX
%================================

\newtcbtheorem[number within=section]{Definition}{Definition}{enhanced,
	before skip=2mm,after skip=2mm, colback=red!5,colframe=red!80!black,boxrule=0.5mm,
	attach boxed title to top left={xshift=1cm,yshift*=1mm-\tcboxedtitleheight}, varwidth boxed title*=-3cm,
	boxed title style={frame code={
					\path[fill=tcbcolback]
					([yshift=-1mm,xshift=-1mm]frame.north west)
					arc[start angle=0,end angle=180,radius=1mm]
					([yshift=-1mm,xshift=1mm]frame.north east)
					arc[start angle=180,end angle=0,radius=1mm];
					\path[left color=tcbcolback!60!black,right color=tcbcolback!60!black,
						middle color=tcbcolback!80!black]
					([xshift=-2mm]frame.north west) -- ([xshift=2mm]frame.north east)
					[rounded corners=1mm]-- ([xshift=1mm,yshift=-1mm]frame.north east)
					-- (frame.south east) -- (frame.south west)
					-- ([xshift=-1mm,yshift=-1mm]frame.north west)
					[sharp corners]-- cycle;
				},interior engine=empty,
		},
	fonttitle=\bfseries,
	title={#2},#1}{def}
\newtcbtheorem[number within=chapter]{definition}{Definition}{enhanced,
	before skip=2mm,after skip=2mm, colback=red!5,colframe=red!80!black,boxrule=0.5mm,
	attach boxed title to top left={xshift=1cm,yshift*=1mm-\tcboxedtitleheight}, varwidth boxed title*=-3cm,
	boxed title style={frame code={
					\path[fill=tcbcolback]
					([yshift=-1mm,xshift=-1mm]frame.north west)
					arc[start angle=0,end angle=180,radius=1mm]
					([yshift=-1mm,xshift=1mm]frame.north east)
					arc[start angle=180,end angle=0,radius=1mm];
					\path[left color=tcbcolback!60!black,right color=tcbcolback!60!black,
						middle color=tcbcolback!80!black]
					([xshift=-2mm]frame.north west) -- ([xshift=2mm]frame.north east)
					[rounded corners=1mm]-- ([xshift=1mm,yshift=-1mm]frame.north east)
					-- (frame.south east) -- (frame.south west)
					-- ([xshift=-1mm,yshift=-1mm]frame.north west)
					[sharp corners]-- cycle;
				},interior engine=empty,
		},
	fonttitle=\bfseries,
	title={#2},#1}{def}



%================================
% Solution BOX
%================================

\makeatletter
\newtcbtheorem{question}{Question}{enhanced,
	breakable,
	colback=white,
	colframe=myb!80!black,
	attach boxed title to top left={yshift*=-\tcboxedtitleheight},
	fonttitle=\bfseries,
	title={#2},
	boxed title size=title,
	boxed title style={%
			sharp corners,
			rounded corners=northwest,
			colback=tcbcolframe,
			boxrule=0pt,
		},
	underlay boxed title={%
			\path[fill=tcbcolframe] (title.south west)--(title.south east)
			to[out=0, in=180] ([xshift=5mm]title.east)--
			(title.center-|frame.east)
			[rounded corners=\kvtcb@arc] |-
			(frame.north) -| cycle;
		},
	#1
}{def}
\makeatother

%================================
% SOLUTION BOX
%================================

\makeatletter
\newtcolorbox{solution}{enhanced,
	breakable,
	colback=white,
	colframe=myg!80!black,
	attach boxed title to top left={yshift*=-\tcboxedtitleheight},
	title=Solution,
	boxed title size=title,
	boxed title style={%
			sharp corners,
			rounded corners=northwest,
			colback=tcbcolframe,
			boxrule=0pt,
		},
	underlay boxed title={%
			\path[fill=tcbcolframe] (title.south west)--(title.south east)
			to[out=0, in=180] ([xshift=5mm]title.east)--
			(title.center-|frame.east)
			[rounded corners=\kvtcb@arc] |-
			(frame.north) -| cycle;
		},
}
\makeatother

%================================
% Question BOX
%================================

\makeatletter
\newtcbtheorem{qstion}{Question}{enhanced,
	breakable,
	colback=white,
	colframe=mygr,
	attach boxed title to top left={yshift*=-\tcboxedtitleheight},
	fonttitle=\bfseries,
	title={#2},
	boxed title size=title,
	boxed title style={%
			sharp corners,
			rounded corners=northwest,
			colback=tcbcolframe,
			boxrule=0pt,
		},
	underlay boxed title={%
			\path[fill=tcbcolframe] (title.south west)--(title.south east)
			to[out=0, in=180] ([xshift=5mm]title.east)--
			(title.center-|frame.east)
			[rounded corners=\kvtcb@arc] |-
			(frame.north) -| cycle;
		},
	#1
}{def}
\makeatother

\newtcbtheorem[number within=chapter]{wconc}{Wrong Concept}{
	breakable,
	enhanced,
	colback=white,
	colframe=myr,
	arc=0pt,
	outer arc=0pt,
	fonttitle=\bfseries\sffamily\large,
	colbacktitle=myr,
	attach boxed title to top left={},
	boxed title style={
			enhanced,
			skin=enhancedfirst jigsaw,
			arc=3pt,
			bottom=0pt,
			interior style={fill=myr}
		},
	#1
}{def}



%================================
% NOTE BOX
%================================

\usetikzlibrary{arrows,calc,shadows.blur}
\tcbuselibrary{skins}
\newtcolorbox{note}[1][]{%
	enhanced jigsaw,
	colback=gray!20!white,%
	colframe=gray!80!black,
	size=small,
	boxrule=1pt,
	title=\textbf{Note:},
	halign title=flush center,
	coltitle=black,
	breakable,
	drop shadow=black!50!white,
	attach boxed title to top left={xshift=1cm,yshift=-\tcboxedtitleheight/2,yshifttext=-\tcboxedtitleheight/2},
	minipage boxed title=1.5cm,
	boxed title style={%
			colback=white,
			size=fbox,
			boxrule=1pt,
			boxsep=2pt,
			underlay={%
					\coordinate (dotA) at ($(interior.west) + (-0.5pt,0)$);
					\coordinate (dotB) at ($(interior.east) + (0.5pt,0)$);
					\begin{scope}
						\clip (interior.north west) rectangle ([xshift=3ex]interior.east);
						\filldraw [white, blur shadow={shadow opacity=60, shadow yshift=-.75ex}, rounded corners=2pt] (interior.north west) rectangle (interior.south east);
					\end{scope}
					\begin{scope}[gray!80!black]
						\fill (dotA) circle (2pt);
						\fill (dotB) circle (2pt);
					\end{scope}
				},
		},
	#1,
}

%%%%%%%%%%%%%%%%%%%%%%%%%%%%%%
% SELF MADE COMMANDS
%%%%%%%%%%%%%%%%%%%%%%%%%%%%%%


\newcommand{\thm}[2]{\begin{Theorem}{#1}{}#2\end{Theorem}}
\newcommand{\cor}[2]{\begin{Corollary}{#1}{}#2\end{Corollary}}
\newcommand{\mlenma}[2]{\begin{Lenma}{#1}{}#2\end{Lenma}}
\newcommand{\mprop}[2]{\begin{Prop}{#1}{}#2\end{Prop}}
\newcommand{\clm}[3]{\begin{claim}{#1}{#2}#3\end{claim}}
\newcommand{\wc}[2]{\begin{wconc}{#1}{}\setlength{\parindent}{1cm}#2\end{wconc}}
\newcommand{\thmcon}[1]{\begin{Theoremcon}{#1}\end{Theoremcon}}
\newcommand{\ex}[2]{\begin{Example}{#1}{}#2\end{Example}}
\newcommand{\dfn}[2]{\begin{Definition}[colbacktitle=red!75!black]{#1}{}#2\end{Definition}}
\newcommand{\dfnc}[2]{\begin{definition}[colbacktitle=red!75!black]{#1}{}#2\end{definition}}
\newcommand{\qs}[2]{\begin{question}{#1}{}#2\end{question}}
\newcommand{\pf}[2]{\begin{myproof}[#1]#2\end{myproof}}
\newcommand{\nt}[1]{\begin{note}#1\end{note}}

\newcommand*\circled[1]{\tikz[baseline=(char.base)]{
		\node[shape=circle,draw,inner sep=1pt] (char) {#1};}}
\newcommand\getcurrentref[1]{%
	\ifnumequal{\value{#1}}{0}
	{??}
	{\the\value{#1}}%
}
\newcommand{\getCurrentSectionNumber}{\getcurrentref{section}}
\newenvironment{myproof}[1][\proofname]{%
	\proof[\bfseries #1: ]%
}{\endproof}

\newcommand{\mclm}[2]{\begin{myclaim}[#1]#2\end{myclaim}}
\newenvironment{myclaim}[1][\claimname]{\proof[\bfseries #1: ]}{}

\newcounter{mylabelcounter}

\makeatletter
\newcommand{\setword}[2]{%
	\phantomsection
	#1\def\@currentlabel{\unexpanded{#1}}\label{#2}%
}
\makeatother




\tikzset{
	symbol/.style={
			draw=none,
			every to/.append style={
					edge node={node [sloped, allow upside down, auto=false]{$#1$}}}
		}
}


% deliminators
\DeclarePairedDelimiter{\abs}{\lvert}{\rvert}
\DeclarePairedDelimiter{\norm}{\lVert}{\rVert}

\DeclarePairedDelimiter{\ceil}{\lceil}{\rceil}
\DeclarePairedDelimiter{\floor}{\lfloor}{\rfloor}
\DeclarePairedDelimiter{\round}{\lfloor}{\rceil}

\newsavebox\diffdbox
\newcommand{\slantedromand}{{\mathpalette\makesl{d}}}
\newcommand{\makesl}[2]{%
\begingroup
\sbox{\diffdbox}{$\mathsurround=0pt#1\mathrm{#2}$}%
\pdfsave
\pdfsetmatrix{1 0 0.2 1}%
\rlap{\usebox{\diffdbox}}%
\pdfrestore
\hskip\wd\diffdbox
\endgroup
}
\newcommand{\dd}[1][]{\ensuremath{\mathop{}\!\ifstrempty{#1}{%
\slantedromand\@ifnextchar^{\hspace{0.2ex}}{\hspace{0.1ex}}}%
{\slantedromand\hspace{0.2ex}^{#1}}}}
\ProvideDocumentCommand\dv{o m g}{%
  \ensuremath{%
    \IfValueTF{#3}{%
      \IfNoValueTF{#1}{%
        \frac{\dd #2}{\dd #3}%
      }{%
        \frac{\dd^{#1} #2}{\dd #3^{#1}}%
      }%
    }{%
      \IfNoValueTF{#1}{%
        \frac{\dd}{\dd #2}%
      }{%
        \frac{\dd^{#1}}{\dd #2^{#1}}%
      }%
    }%
  }%
}
\providecommand*{\pdv}[3][]{\frac{\partial^{#1}#2}{\partial#3^{#1}}}
%  - others
\DeclareMathOperator{\Lap}{\mathcal{L}}
\DeclareMathOperator{\Var}{Var} % varience
\DeclareMathOperator{\Cov}{Cov} % covarience
\DeclareMathOperator{\E}{E} % expected

% Since the amsthm package isn't loaded

% I prefer the slanted \leq
\let\oldleq\leq % save them in case they're every wanted
\let\oldgeq\geq
\renewcommand{\leq}{\leqslant}
\renewcommand{\geq}{\geqslant}

% % redefine matrix env to allow for alignment, use r as default
% \renewcommand*\env@matrix[1][r]{\hskip -\arraycolsep
%     \let\@ifnextchar\new@ifnextchar
%     \array{*\c@MaxMatrixCols #1}}


%\usepackage{framed}
%\usepackage{titletoc}
%\usepackage{etoolbox}
%\usepackage{lmodern}


%\patchcmd{\tableofcontents}{\contentsname}{\sffamily\contentsname}{}{}

%\renewenvironment{leftbar}
%{\def\FrameCommand{\hspace{6em}%
%		{\color{myyellow}\vrule width 2pt depth 6pt}\hspace{1em}}%
%	\MakeFramed{\parshape 1 0cm \dimexpr\textwidth-6em\relax\FrameRestore}\vskip2pt%
%}
%{\endMakeFramed}

%\titlecontents{chapter}
%[0em]{\vspace*{2\baselineskip}}
%{\parbox{4.5em}{%
%		\hfill\Huge\sffamily\bfseries\color{myred}\thecontentspage}%
%	\vspace*{-2.3\baselineskip}\leftbar\textsc{\small\chaptername~\thecontentslabel}\\\sffamily}
%{}{\endleftbar}
%\titlecontents{section}
%[8.4em]
%{\sffamily\contentslabel{3em}}{}{}
%{\hspace{0.5em}\nobreak\itshape\color{myred}\contentspage}
%\titlecontents{subsection}
%[8.4em]
%{\sffamily\contentslabel{3em}}{}{}  
%{\hspace{0.5em}\nobreak\itshape\color{myred}\contentspage}



%%%%%%%%%%%%%%%%%%%%%%%%%%%%%%%%%%%%%%%%%%%
% TABLE OF CONTENTS
%%%%%%%%%%%%%%%%%%%%%%%%%%%%%%%%%%%%%%%%%%%

\usepackage{tikz}
\definecolor{doc}{RGB}{0,60,110}
\usepackage{titletoc}
\contentsmargin{0cm}
\titlecontents{chapter}[3.7pc]
{\addvspace{30pt}%
	\begin{tikzpicture}[remember picture, overlay]%
		\draw[fill=doc!60,draw=doc!60] (-7,-.1) rectangle (-0.9,.5);%
		\pgftext[left,x=-3.5cm,y=0.2cm]{\color{white}\Large\sc\bfseries Chapter\ \thecontentslabel};%
	\end{tikzpicture}\color{doc!60}\large\sc\bfseries}%
{}
{}
{\;\titlerule\;\large\sc\bfseries Page \thecontentspage
	\begin{tikzpicture}[remember picture, overlay]
		\draw[fill=doc!60,draw=doc!60] (2pt,0) rectangle (4,0.1pt);
	\end{tikzpicture}}%
\titlecontents{section}[3.7pc]
{\addvspace{2pt}}
{\contentslabel[\thecontentslabel]{2pc}}
{}
{\hfill\small \thecontentspage}
[]
\titlecontents*{subsection}[3.7pc]
{\addvspace{-1pt}\small}
{}
{}
{\ --- \small\thecontentspage}
[ \textbullet\ ][]

\makeatletter
\renewcommand{\tableofcontents}{%
	\chapter*{%
	  \vspace*{-20\p@}%
	  \begin{tikzpicture}[remember picture, overlay]%
		  \pgftext[right,x=15cm,y=0.2cm]{\color{doc!60}\Huge\sc\bfseries \contentsname};%
		  \draw[fill=doc!60,draw=doc!60] (13,-.75) rectangle (20,1);%
		  \clip (13,-.75) rectangle (20,1);
		  \pgftext[right,x=15cm,y=0.2cm]{\color{white}\Huge\sc\bfseries \contentsname};%
	  \end{tikzpicture}}%
	\@starttoc{toc}}
\makeatother


% Comandi di Bastiality %
\newcommand{\innerproduct}[2]{\langle #1, #2 \rangle}
\newcommand{\generators}[1]{\langle #1 \rangle}
\newcommand{\powerset}[0]{\mathcal{P}}
\newcommand{\indip}{\perp \!\!\! \perp}

% Comandi di GioLaPalma %
\newcommand{\der}[0]{\implies^*}
\newcommand{\red}[1]{\textcolor{red}{#1}}
\newcommand{\esempio}[1]{\ex{}{#1}}
\newcommand{\teorema}[1]{\thm{}{#1}}
\newcommand{\dimostrazione}[1]{\pf{Dimostrazione}{#1}}
\setlength{\parindent}{0pt}
\newcommand{\osservazione}[1]{\nt{#1}}
\newcommand{\inside}[1]{
  \langle #1 \rangle
}
      

%From M275 "Topology" at SJSU
\newcommand{\id}{\mathrm{id}}
\newcommand{\taking}[1]{\xrightarrow{#1}}
\newcommand{\inv}{^{-1}}

%From M170 "Introduction to Graph Theory" at SJSU
\DeclareMathOperator{\diam}{diam}
\DeclareMathOperator{\ord}{ord}
\newcommand{\defeq}{\overset{\mathrm{def}}{=}}

%From the USAMO .tex files
\newcommand{\ts}{\textsuperscript}
\newcommand{\dg}{^\circ}
\newcommand{\ii}{\item}

% % From Math 55 and Math 145 at Harvard
% \newenvironment{subproof}[1][Proof]{%
% \begin{proof}[#1] \renewcommand{\qedsymbol}{$\blacksquare$}}%
% {\end{proof}}

\newcommand{\liff}{\leftrightarrow}
\newcommand{\lthen}{\rightarrow}
\newcommand{\opname}{\operatorname}
\newcommand{\surjto}{\twoheadrightarrow}
\newcommand{\injto}{\hookrightarrow}
\newcommand{\On}{\mathrm{On}} % ordinals
\DeclareMathOperator{\img}{im} % Image
\DeclareMathOperator{\Img}{Im} % Image
\DeclareMathOperator{\coker}{coker} % Cokernel
\DeclareMathOperator{\Coker}{Coker} % Cokernel
\DeclareMathOperator{\Ker}{Ker} % Kernel
\DeclareMathOperator{\rank}{rank}
\DeclareMathOperator{\Spec}{Spec} % spectrum
\DeclareMathOperator{\Tr}{Tr} % trace
\DeclareMathOperator{\pr}{pr} % projection
\DeclareMathOperator{\ext}{ext} % extension
\DeclareMathOperator{\pred}{pred} % predecessor
\DeclareMathOperator{\dom}{dom} % domain
\DeclareMathOperator{\ran}{ran} % range
\DeclareMathOperator{\Hom}{Hom} % homomorphism
\DeclareMathOperator{\Mor}{Mor} % morphisms
\DeclareMathOperator{\End}{End} % endomorphism

\newcommand{\eps}{\epsilon}
\newcommand{\veps}{\varepsilon}
\newcommand{\ol}{\overline}
\newcommand{\ul}{\underline}
\newcommand{\wt}{\widetilde}
\newcommand{\wh}{\widehat}
\newcommand{\vocab}[1]{\textbf{\color{blue} #1}}
\providecommand{\half}{\frac{1}{2}}
\newcommand{\dang}{\measuredangle} %% Directed angle
\newcommand{\ray}[1]{\overrightarrow{#1}}
\newcommand{\seg}[1]{\overline{#1}}
\newcommand{\arc}[1]{\wideparen{#1}}
\DeclareMathOperator{\cis}{cis}
\DeclareMathOperator*{\lcm}{lcm}
\DeclareMathOperator*{\argmin}{arg min}
\DeclareMathOperator*{\argmax}{arg max}
\newcommand{\cycsum}{\sum_{\mathrm{cyc}}}
\newcommand{\symsum}{\sum_{\mathrm{sym}}}
\newcommand{\cycprod}{\prod_{\mathrm{cyc}}}
\newcommand{\symprod}{\prod_{\mathrm{sym}}}
\newcommand{\Qed}{\begin{flushright}\qed\end{flushright}}
\newcommand{\parinn}{\setlength{\parindent}{1cm}}
\newcommand{\parinf}{\setlength{\parindent}{0cm}}
% \newcommand{\norm}{\|\cdot\|}
\newcommand{\inorm}{\norm_{\infty}}
\newcommand{\opensets}{\{V_{\alpha}\}_{\alpha\in I}}
\newcommand{\oset}{V_{\alpha}}
\newcommand{\opset}[1]{V_{\alpha_{#1}}}
\newcommand{\lub}{\text{lub}}
\newcommand{\del}[2]{\frac{\partial #1}{\partial #2}}
\newcommand{\Del}[3]{\frac{\partial^{#1} #2}{\partial^{#1} #3}}
\newcommand{\deld}[2]{\dfrac{\partial #1}{\partial #2}}
\newcommand{\Deld}[3]{\dfrac{\partial^{#1} #2}{\partial^{#1} #3}}
\newcommand{\lm}{\lambda}
\newcommand{\uin}{\mathbin{\rotatebox[origin=c]{90}{$\in$}}}
\newcommand{\usubset}{\mathbin{\rotatebox[origin=c]{90}{$\subset$}}}
\newcommand{\lt}{\left}
\newcommand{\rt}{\right}
\newcommand{\bs}[1]{\boldsymbol{#1}}
\newcommand{\exs}{\exists}
\newcommand{\st}{\strut}
\newcommand{\dps}[1]{\displaystyle{#1}}

\newcommand{\sol}{\setlength{\parindent}{0cm}\textbf{\textit{Solution:}}\setlength{\parindent}{1cm} }
\newcommand{\solve}[1]{\setlength{\parindent}{0cm}\textbf{\textit{Solution: }}\setlength{\parindent}{1cm}#1 \Qed}

% Things Lie
\newcommand{\kb}{\mathfrak b}
\newcommand{\kg}{\mathfrak g}
\newcommand{\kh}{\mathfrak h}
\newcommand{\kn}{\mathfrak n}
\newcommand{\ku}{\mathfrak u}
\newcommand{\kz}{\mathfrak z}
\DeclareMathOperator{\Ext}{Ext} % Ext functor
\DeclareMathOperator{\Tor}{Tor} % Tor functor
\newcommand{\gl}{\opname{\mathfrak{gl}}} % frak gl group
\renewcommand{\sl}{\opname{\mathfrak{sl}}} % frak sl group chktex 6

% More script letters etc.
\newcommand{\SA}{\mathcal A}
\newcommand{\SB}{\mathcal B}
\newcommand{\SC}{\mathcal C}
\newcommand{\SF}{\mathcal F}
\newcommand{\SG}{\mathcal G}
\newcommand{\SH}{\mathcal H}
\newcommand{\OO}{\mathcal O}

\newcommand{\SCA}{\mathscr A}
\newcommand{\SCB}{\mathscr B}
\newcommand{\SCC}{\mathscr C}
\newcommand{\SCD}{\mathscr D}
\newcommand{\SCE}{\mathscr E}
\newcommand{\SCF}{\mathscr F}
\newcommand{\SCG}{\mathscr G}
\newcommand{\SCH}{\mathscr H}

% Mathfrak primes
\newcommand{\km}{\mathfrak m}
\newcommand{\kp}{\mathfrak p}
\newcommand{\kq}{\mathfrak q}

% number sets
\newcommand{\RR}[1][]{\ensuremath{\ifstrempty{#1}{\mathbb{R}}{\mathbb{R}^{#1}}}}
\newcommand{\NN}[1][]{\ensuremath{\ifstrempty{#1}{\mathbb{N}}{\mathbb{N}^{#1}}}}
\newcommand{\ZZ}[1][]{\ensuremath{\ifstrempty{#1}{\mathbb{Z}}{\mathbb{Z}^{#1}}}}
\newcommand{\QQ}[1][]{\ensuremath{\ifstrempty{#1}{\mathbb{Q}}{\mathbb{Q}^{#1}}}}
\newcommand{\CC}[1][]{\ensuremath{\ifstrempty{#1}{\mathbb{C}}{\mathbb{C}^{#1}}}}
\newcommand{\PP}[1][]{\ensuremath{\ifstrempty{#1}{\mathbb{P}}{\mathbb{P}^{#1}}}}
\newcommand{\HH}[1][]{\ensuremath{\ifstrempty{#1}{\mathbb{H}}{\mathbb{H}^{#1}}}}
\newcommand{\FF}[1][]{\ensuremath{\ifstrempty{#1}{\mathbb{F}}{\mathbb{F}^{#1}}}}
% expected value
\newcommand{\EE}{\ensuremath{\mathbb{E}}}
\newcommand{\charin}{\text{ char }}
\DeclareMathOperator{\sign}{sign}
\DeclareMathOperator{\Aut}{Aut}
\DeclareMathOperator{\Inn}{Inn}
\DeclareMathOperator{\Syl}{Syl}
\DeclareMathOperator{\Gal}{Gal}
\DeclareMathOperator{\GL}{GL} % General linear group
\DeclareMathOperator{\SL}{SL} % Special linear group

%---------------------------------------
% BlackBoard Math Fonts :-
%---------------------------------------

%Captital Letters
\newcommand{\bbA}{\mathbb{A}}	\newcommand{\bbB}{\mathbb{B}}
\newcommand{\bbC}{\mathbb{C}}	\newcommand{\bbD}{\mathbb{D}}
\newcommand{\bbE}{\mathbb{E}}	\newcommand{\bbF}{\mathbb{F}}
\newcommand{\bbG}{\mathbb{G}}	\newcommand{\bbH}{\mathbb{H}}
\newcommand{\bbI}{\mathbb{I}}	\newcommand{\bbJ}{\mathbb{J}}
\newcommand{\bbK}{\mathbb{K}}	\newcommand{\bbL}{\mathbb{L}}
\newcommand{\bbM}{\mathbb{M}}	\newcommand{\bbN}{\mathbb{N}}
\newcommand{\bbO}{\mathbb{O}}	\newcommand{\bbP}{\mathbb{P}}
\newcommand{\bbQ}{\mathbb{Q}}	\newcommand{\bbR}{\mathbb{R}}
\newcommand{\bbS}{\mathbb{S}}	\newcommand{\bbT}{\mathbb{T}}
\newcommand{\bbU}{\mathbb{U}}	\newcommand{\bbV}{\mathbb{V}}
\newcommand{\bbW}{\mathbb{W}}	\newcommand{\bbX}{\mathbb{X}}
\newcommand{\bbY}{\mathbb{Y}}	\newcommand{\bbZ}{\mathbb{Z}}

%---------------------------------------
% MathCal Fonts :-
%---------------------------------------

%Captital Letters
\newcommand{\mcA}{\mathcal{A}}	\newcommand{\mcB}{\mathcal{B}}
\newcommand{\mcC}{\mathcal{C}}	\newcommand{\mcD}{\mathcal{D}}
\newcommand{\mcE}{\mathcal{E}}	\newcommand{\mcF}{\mathcal{F}}
\newcommand{\mcG}{\mathcal{G}}	\newcommand{\mcH}{\mathcal{H}}
\newcommand{\mcI}{\mathcal{I}}	\newcommand{\mcJ}{\mathcal{J}}
\newcommand{\mcK}{\mathcal{K}}	\newcommand{\mcL}{\mathcal{L}}
\newcommand{\mcM}{\mathcal{M}}	\newcommand{\mcN}{\mathcal{N}}
\newcommand{\mcO}{\mathcal{O}}	\newcommand{\mcP}{\mathcal{P}}
\newcommand{\mcQ}{\mathcal{Q}}	\newcommand{\mcR}{\mathcal{R}}
\newcommand{\mcS}{\mathcal{S}}	\newcommand{\mcT}{\mathcal{T}}
\newcommand{\mcU}{\mathcal{U}}	\newcommand{\mcV}{\mathcal{V}}
\newcommand{\mcW}{\mathcal{W}}	\newcommand{\mcX}{\mathcal{X}}
\newcommand{\mcY}{\mathcal{Y}}	\newcommand{\mcZ}{\mathcal{Z}}


%---------------------------------------
% Bold Math Fonts :-
%---------------------------------------

%Captital Letters
\newcommand{\bmA}{\boldsymbol{A}}	\newcommand{\bmB}{\boldsymbol{B}}
\newcommand{\bmC}{\boldsymbol{C}}	\newcommand{\bmD}{\boldsymbol{D}}
\newcommand{\bmE}{\boldsymbol{E}}	\newcommand{\bmF}{\boldsymbol{F}}
\newcommand{\bmG}{\boldsymbol{G}}	\newcommand{\bmH}{\boldsymbol{H}}
\newcommand{\bmI}{\boldsymbol{I}}	\newcommand{\bmJ}{\boldsymbol{J}}
\newcommand{\bmK}{\boldsymbol{K}}	\newcommand{\bmL}{\boldsymbol{L}}
\newcommand{\bmM}{\boldsymbol{M}}	\newcommand{\bmN}{\boldsymbol{N}}
\newcommand{\bmO}{\boldsymbol{O}}	\newcommand{\bmP}{\boldsymbol{P}}
\newcommand{\bmQ}{\boldsymbol{Q}}	\newcommand{\bmR}{\boldsymbol{R}}
\newcommand{\bmS}{\boldsymbol{S}}	\newcommand{\bmT}{\boldsymbol{T}}
\newcommand{\bmU}{\boldsymbol{U}}	\newcommand{\bmV}{\boldsymbol{V}}
\newcommand{\bmW}{\boldsymbol{W}}	\newcommand{\bmX}{\boldsymbol{X}}
\newcommand{\bmY}{\boldsymbol{Y}}	\newcommand{\bmZ}{\boldsymbol{Z}}
%Small Letters
\newcommand{\bma}{\boldsymbol{a}}	\newcommand{\bmb}{\boldsymbol{b}}
\newcommand{\bmc}{\boldsymbol{c}}	\newcommand{\bmd}{\boldsymbol{d}}
\newcommand{\bme}{\boldsymbol{e}}	\newcommand{\bmf}{\boldsymbol{f}}
\newcommand{\bmg}{\boldsymbol{g}}	\newcommand{\bmh}{\boldsymbol{h}}
\newcommand{\bmi}{\boldsymbol{i}}	\newcommand{\bmj}{\boldsymbol{j}}
\newcommand{\bmk}{\boldsymbol{k}}	\newcommand{\bml}{\boldsymbol{l}}
\newcommand{\bmm}{\boldsymbol{m}}	\newcommand{\bmn}{\boldsymbol{n}}
\newcommand{\bmo}{\boldsymbol{o}}	\newcommand{\bmp}{\boldsymbol{p}}
\newcommand{\bmq}{\boldsymbol{q}}	\newcommand{\bmr}{\boldsymbol{r}}
\newcommand{\bms}{\boldsymbol{s}}	\newcommand{\bmt}{\boldsymbol{t}}
\newcommand{\bmu}{\boldsymbol{u}}	\newcommand{\bmv}{\boldsymbol{v}}
\newcommand{\bmw}{\boldsymbol{w}}	\newcommand{\bmx}{\boldsymbol{x}}
\newcommand{\bmy}{\boldsymbol{y}}	\newcommand{\bmz}{\boldsymbol{z}}

%---------------------------------------
% Scr Math Fonts :-
%---------------------------------------

\newcommand{\sA}{{\mathscr{A}}}   \newcommand{\sB}{{\mathscr{B}}}
\newcommand{\sC}{{\mathscr{C}}}   \newcommand{\sD}{{\mathscr{D}}}
\newcommand{\sE}{{\mathscr{E}}}   \newcommand{\sF}{{\mathscr{F}}}
\newcommand{\sG}{{\mathscr{G}}}   \newcommand{\sH}{{\mathscr{H}}}
\newcommand{\sI}{{\mathscr{I}}}   \newcommand{\sJ}{{\mathscr{J}}}
\newcommand{\sK}{{\mathscr{K}}}   \newcommand{\sL}{{\mathscr{L}}}
\newcommand{\sM}{{\mathscr{M}}}   \newcommand{\sN}{{\mathscr{N}}}
\newcommand{\sO}{{\mathscr{O}}}   \newcommand{\sP}{{\mathscr{P}}}
\newcommand{\sQ}{{\mathscr{Q}}}   \newcommand{\sR}{{\mathscr{R}}}
\newcommand{\sS}{{\mathscr{S}}}   \newcommand{\sT}{{\mathscr{T}}}
\newcommand{\sU}{{\mathscr{U}}}   \newcommand{\sV}{{\mathscr{V}}}
\newcommand{\sW}{{\mathscr{W}}}   \newcommand{\sX}{{\mathscr{X}}}
\newcommand{\sY}{{\mathscr{Y}}}   \newcommand{\sZ}{{\mathscr{Z}}}


%---------------------------------------
% Math Fraktur Font
%---------------------------------------

%Captital Letters
\newcommand{\mfA}{\mathfrak{A}}	\newcommand{\mfB}{\mathfrak{B}}
\newcommand{\mfC}{\mathfrak{C}}	\newcommand{\mfD}{\mathfrak{D}}
\newcommand{\mfE}{\mathfrak{E}}	\newcommand{\mfF}{\mathfrak{F}}
\newcommand{\mfG}{\mathfrak{G}}	\newcommand{\mfH}{\mathfrak{H}}
\newcommand{\mfI}{\mathfrak{I}}	\newcommand{\mfJ}{\mathfrak{J}}
\newcommand{\mfK}{\mathfrak{K}}	\newcommand{\mfL}{\mathfrak{L}}
\newcommand{\mfM}{\mathfrak{M}}	\newcommand{\mfN}{\mathfrak{N}}
\newcommand{\mfO}{\mathfrak{O}}	\newcommand{\mfP}{\mathfrak{P}}
\newcommand{\mfQ}{\mathfrak{Q}}	\newcommand{\mfR}{\mathfrak{R}}
\newcommand{\mfS}{\mathfrak{S}}	\newcommand{\mfT}{\mathfrak{T}}
\newcommand{\mfU}{\mathfrak{U}}	\newcommand{\mfV}{\mathfrak{V}}
\newcommand{\mfW}{\mathfrak{W}}	\newcommand{\mfX}{\mathfrak{X}}
\newcommand{\mfY}{\mathfrak{Y}}	\newcommand{\mfZ}{\mathfrak{Z}}
%Small Letters
\newcommand{\mfa}{\mathfrak{a}}	\newcommand{\mfb}{\mathfrak{b}}
\newcommand{\mfc}{\mathfrak{c}}	\newcommand{\mfd}{\mathfrak{d}}
\newcommand{\mfe}{\mathfrak{e}}	\newcommand{\mff}{\mathfrak{f}}
\newcommand{\mfg}{\mathfrak{g}}	\newcommand{\mfh}{\mathfrak{h}}
\newcommand{\mfi}{\mathfrak{i}}	\newcommand{\mfj}{\mathfrak{j}}
\newcommand{\mfk}{\mathfrak{k}}	\newcommand{\mfl}{\mathfrak{l}}
\newcommand{\mfm}{\mathfrak{m}}	\newcommand{\mfn}{\mathfrak{n}}
\newcommand{\mfo}{\mathfrak{o}}	\newcommand{\mfp}{\mathfrak{p}}
\newcommand{\mfq}{\mathfrak{q}}	\newcommand{\mfr}{\mathfrak{r}}
\newcommand{\mfs}{\mathfrak{s}}	\newcommand{\mft}{\mathfrak{t}}
\newcommand{\mfu}{\mathfrak{u}}	\newcommand{\mfv}{\mathfrak{v}}
\newcommand{\mfw}{\mathfrak{w}}	\newcommand{\mfx}{\mathfrak{x}}
\newcommand{\mfy}{\mathfrak{y}}	\newcommand{\mfz}{\mathfrak{z}}


\usepackage[T1]{fontenc}
\usepackage[utf8]{inputenc}
\usepackage{xcolor}
\usepackage{graphicx}
\usepackage{listings}
\usepackage{array}
\usepackage{longtable}

%==============================================================
% Stile per i blocchi di codice (C, C++, Java, Python)
%==============================================================
\definecolor{codebg}{HTML}{F7F7FB}
\definecolor{codekw}{HTML}{2D6FB1}
\definecolor{codecm}{HTML}{388C46}
\definecolor{codestr}{HTML}{C74440}
\definecolor{coderule}{HTML}{CCCCCC}

\lstset{
    basicstyle=\ttfamily\footnotesize,
    keywordstyle=\color{codekw}\bfseries,
    commentstyle=\color{codecm}\itshape,
    stringstyle=\color{codestr},
    numberstyle=\tiny\color{gray},
    showstringspaces=false,
    breaklines=true,
    breakatwhitespace=false,
    frame=single,
    rulecolor=\color{coderule},
    backgroundcolor=\color{codebg},
    tabsize=2,
    columns=fullflexible,
    keepspaces=true,
    captionpos=b,
    xleftmargin=1.2em,
    aboveskip=1em,
    belowskip=1em,
}

% codice inline
\newcommand{\code}[1]{\texttt{\small #1}}
% riga di intestazione linguaggio dentro un listing
\newcommand{\lang}[1]{\hfill\textbf{\footnotesize\color{codekw}#1}}

\graphicspath{{./img/}}

%==============================================================
\begin{document}

\begin{titlepage}
    \centering
    \vspace*{3cm}
    {\Huge\sffamily\bfseries Differenze tra i linguaggi\par}
    \vspace{0.6cm}
    {\LARGE\sffamily C \; $\cdot$ \; C++ \; $\cdot$ \; Java \; $\cdot$ \; Python\par}
    \vspace{1.2cm}
    {\Large Appunti di Ingegneria del Software\par}
    \vspace{0.4cm}
    {\large Funzioni virtuali \& late binding \; --- \; Template \& generics \; --- \; Strategy pattern\par}
    \vspace{2cm}
    {\large\itshape Rielaborazione trasversale delle slide del corso\par}
    \vfill
    {\large Anno 3\par}
\end{titlepage}

\tableofcontents
\newpage

%==============================================================
\chapter{Introduzione: perché confrontare i linguaggi}
%==============================================================

Le slide del corso presentano tre argomenti (\textbf{funzioni virtuali e late
binding}, \textbf{template}, \textbf{strategy pattern}) usando principalmente
C++ e Java come linguaggi di riferimento. Il professore, però, chiede
ripetutamente di ragionare sulle \emph{differenze tra i linguaggi}: perché una
certa feature esiste in C++ ma non in Java, come si emula un meccanismo in C
che non ha l'orientamento agli oggetti, cosa cambia in un linguaggio a
tipizzazione dinamica come Python.

Questi appunti riorganizzano quindi \emph{tutti} gli argomenti delle slide
mettendo sistematicamente a confronto \textbf{C, C++, Java e Python}. L'idea è
che ogni meccanismo (l'overriding, il template, lo strategy) va capito non
come un dettaglio sintattico di un linguaggio, ma come \emph{una risposta a un
problema di ingegneria del software} (riuso, estensibilità, disaccoppiamento)
che ciascun linguaggio risolve con scelte diverse.

\section{I quattro linguaggi in una pagina}

Prima di entrare nei dettagli conviene fissare le coordinate fondamentali dei
quattro linguaggi, perché quasi tutte le differenze che vedremo discendono da
queste scelte di progetto.

\dfn{Tipizzazione statica vs dinamica}{
    Un linguaggio è a \textbf{tipizzazione statica} se il tipo di ogni
    espressione è noto e controllato a \emph{tempo di compilazione} (C, C++,
    Java). È a \textbf{tipizzazione dinamica} se i tipi sono associati ai
    \emph{valori} e controllati a \emph{tempo di esecuzione} (Python).
    \\[2pt]
    Ortogonale a questa c'è la distinzione \textbf{forte/debole}: un tipaggio è
    \emph{forte} se il linguaggio impedisce (o rende esplicite) le conversioni
    di tipo pericolose. C è statico ma relativamente \emph{debole} (i cast e
    l'aritmetica dei puntatori permettono di reinterpretare la memoria); Java e
    Python sono \emph{forti}.
}

\begin{center}
\renewcommand{\arraystretch}{1.35}
\footnotesize
\begin{tabular}{|>{\raggedright\arraybackslash}p{2.6cm}|>{\raggedright\arraybackslash}p{2.5cm}|>{\raggedright\arraybackslash}p{2.8cm}|>{\raggedright\arraybackslash}p{2.6cm}|>{\raggedright\arraybackslash}p{2.6cm}|}
\hline
\textbf{Aspetto} & \textbf{C} & \textbf{C++} & \textbf{Java} & \textbf{Python} \\
\hline
Paradigma & imperativo / procedurale & multiparadigma (proc. + OOP + generico) & OOP (class-based) & multiparadigma (OOP + funzionale) \\
\hline
Tipizzazione & statica, debole & statica, forte & statica, forte & dinamica, forte \\
\hline
Esecuzione & compilato $\to$ nativo & compilato $\to$ nativo & compilato $\to$ bytecode $\to$ JVM & interpretato $\to$ bytecode $\to$ VM (CPython) \\
\hline
Memoria & manuale (\code{malloc}/\code{free}) & manuale + RAII / smart pointer & garbage collector & GC (reference counting + cicli) \\
\hline
Dispatch metodi & assente (puntatori a funzione) & statico di default, dinamico se \code{virtual} & \textbf{sempre} dinamico & \textbf{sempre} dinamico (duck typing) \\
\hline
Genericità & macro, \code{void*}, \code{\_Generic} & template (monomorfizzazione) & generics (type erasure) & duck typing + type hints \\
\hline
Ereditarietà & assente & multipla (con \code{virtual}) & singola + interfacce multiple & multipla (MRO C3) \\
\hline
\end{tabular}
\end{center}

\nt{
    Tieni a mente questa tabella: è la \emph{chiave di lettura} di tutto il
    resto. Ogni volta che una feature "esiste in un linguaggio e non
    nell'altro" la spiegazione sta quasi sempre in una di queste righe
    (tipizzazione, modello di esecuzione, gestione della memoria).
}

\section{Modello di esecuzione: cosa succede al codice}

\begin{itemize}
    \item \textbf{C e C++} sono \emph{compilati a codice macchina nativo}. Il
    compilatore traduce direttamente in istruzioni della CPU; il binario gira
    "sul metallo". Questo dà massime prestazioni e controllo, ma nessuna rete
    di sicurezza a runtime.
    \item \textbf{Java} è compilato in un formato intermedio (il
    \emph{bytecode}) eseguito da una macchina virtuale (JVM). La JVM aggiunge
    controlli a runtime (array bounds, type safety, garbage collection) e
    porta il compromesso "write once, run anywhere".
    \item \textbf{Python} (CPython) compila il sorgente in bytecode che poi
    \emph{interpreta} in un ciclo di valutazione. Nessun controllo statico dei
    tipi: gli errori di tipo emergono solo quando la riga viene eseguita.
\end{itemize}

Questa gerarchia (nativo $\to$ bytecode+VM $\to$ interpretato) spiega il
trade-off ricorrente del corso: \textbf{più controllo e prestazioni} da un
lato (C/C++), \textbf{più sicurezza e produttività} dall'altro (Java/Python).

%==============================================================
\chapter{Binding e funzioni virtuali}
%==============================================================

Questo capitolo copre la lezione \emph{Virtual Functions and Late Bindings}. Il
concetto centrale è il \textbf{binding}: il momento in cui una chiamata di
funzione viene associata al codice che effettivamente verrà eseguito.

\section{Early binding vs late binding}

\dfn{Binding}{
    Il \textbf{binding} è l'associazione tra una \emph{chiamata} a funzione (il
    punto nel codice dove appare l'invocazione) e la \emph{definizione} della
    funzione (l'indirizzo di memoria del codice da eseguire).
}

\dfn{Early binding (compile-time binding)}{
    L'associazione chiamata$\to$codice avviene a \textbf{tempo di
    compilazione}. Il compilatore, incontrando la chiamata, scrive
    un'istruzione macchina di \code{call} verso un indirizzo ben preciso e
    fisso. È il comportamento standard delle funzioni in C e delle funzioni
    \emph{non} virtuali in C++.
}

\dfn{Late binding (run-time binding)}{
    L'associazione chiamata$\to$codice viene rimandata al \textbf{tempo di
    esecuzione}: il compilatore non genera una \code{call} a un indirizzo
    fisso, ma codice che \emph{calcola} a runtime quale funzione invocare, in
    base al tipo \emph{effettivo} (dinamico) dell'oggetto. È il meccanismo che
    rende possibili le funzioni virtuali.
}

\nt{
    Uno dei principi centrali dell'OOP è determinare il comportamento degli
    oggetti \emph{il più tardi possibile}, idealmente a runtime. Nei linguaggi
    imperativi non-OO (C, Pascal, \dots) l'unico modo per farlo sono \code{if}
    o \code{switch} espliciti sul "tipo": le funzioni virtuali sono la
    soluzione elegante e automatica a questo problema.
}

\subsection{Il confronto fondamentale sul dispatch}

\begin{center}
\renewcommand{\arraystretch}{1.35}
\footnotesize
\begin{tabular}{|>{\raggedright\arraybackslash}p{2.3cm}|>{\raggedright\arraybackslash}p{9.5cm}|}
\hline
\textbf{Linguaggio} & \textbf{Comportamento di default del dispatch} \\
\hline
C & Nessun dispatch dinamico nativo. Le chiamate sono \emph{sempre} in early
binding. Il polimorfismo si emula a mano con \textbf{puntatori a funzione}
(spesso raccolti in una \code{struct}, che diventa una vtable manuale). \\
\hline
C++ & Early binding di \textbf{default}. Il late binding si attiva solo
dichiarando il metodo \code{virtual}. Filosofia: "non paghi per ciò che non
usi" (\emph{zero-overhead}). \\
\hline
Java & Late binding \textbf{sempre}: tutti i metodi non statici sono virtuali
per default. Si può impedire l'override con \code{final}. \\
\hline
Python & Late binding \textbf{sempre}, per natura del linguaggio dinamico: il
metodo da chiamare è cercato a runtime nel dizionario della classe (duck
typing). Non esiste il concetto di "virtuale/non virtuale". \\
\hline
\end{tabular}
\end{center}

\section{Funzioni virtuali e overriding}

\dfn{Funzione virtuale}{
    Un metodo \emph{non statico} di una classe base che può essere
    \textbf{ridefinito} (\emph{overridden}) in una classe derivata. La
    ridefinizione fa sì che, chiamando il metodo su un riferimento/puntatore
    alla classe base, venga eseguita la versione della classe \emph{reale}
    dell'oggetto.
}

\dfn{Overriding}{
    Il processo con cui un metodo della classe base viene \emph{sostituito} da
    un altro nella classe derivata. Avviene solo se:
    \begin{itemize}
        \item il metodo sovrascritto è virtuale;
        \item il metodo nella derivata ha \textbf{stessa signature} (e stesso
        tipo di ritorno) del metodo nella base.
    \end{itemize}
}

\dfn{Signature e prototipo}{
    La \textbf{signature} di una funzione comprende: \emph{nome}, \emph{numero}
    dei parametri formali, \emph{tipi} dei parametri formali. Il \emph{tipo di
    ritorno NON fa parte della signature}. \\[2pt]
    La signature \emph{insieme} al tipo di ritorno è detta \textbf{prototipo}
    della funzione.
}

\mprop{Le tre possibilità dell'invocazione virtuale}{
    Quando una funzione virtuale $f$ è invocata su un oggetto $o$ di classe $D$
    derivata da $B$, ci sono tre casi:
    \begin{enumerate}
        \item $f$ definita solo in $D$: viene invocata $D.f$;
        \item $f$ definita solo in $B$: viene invocata $B.f$ (ereditata da $D$);
        \item $f$ definita (overridden) sia in $B$ sia in $D$: viene invocata
        $D.f$.
    \end{enumerate}
    In tutti i casi conta la \textbf{classe reale} dell'oggetto, non la classe
    del riferimento che lo punta.
}

\nt{
    Quando si invoca una funzione virtuale su un oggetto il cui tipo è
    \emph{noto a compile-time} (es. un oggetto locale, non un
    riferimento/puntatore), il comportamento è indistinguibile da quello di una
    funzione non virtuale: il vantaggio del polimorfismo emerge \emph{solo}
    manipolando gli oggetti tramite riferimenti (Java) o puntatori (C++) alla
    classe base.
}

\subsection{Overriding vs hiding vs overloading}

Un punto su cui il prof insiste e che distingue i linguaggi:

\begin{itemize}
    \item \textbf{Overriding}: stessa signature + metodo virtuale $\Rightarrow$
    sostituzione a runtime (polimorfismo).
    \item \textbf{Overloading}: stesso nome, signature \emph{diversa}
    $\Rightarrow$ funzioni distinte scelte a compile-time in base agli
    argomenti.
    \item \textbf{Hiding}: in C++, se nella derivata dichiaro un metodo con lo
    stesso nome ma signature diversa di uno della base, \emph{tutti} i metodi
    con quel nome nella base diventano \textbf{invisibili} nella derivata (non
    c'è overloading tra base e derivata).
\end{itemize}

\begin{center}
\renewcommand{\arraystretch}{1.35}
\footnotesize
\begin{tabular}{|>{\raggedright\arraybackslash}p{3.2cm}|>{\raggedright\arraybackslash}p{4.2cm}|>{\raggedright\arraybackslash}p{4.6cm}|}
\hline
\textbf{Situazione} & \textbf{C++} & \textbf{Java} \\
\hline
Metodo non dichiarato virtuale, ridefinito nella derivata &
Non è override: è \textbf{hiding} (early binding sul tipo statico). &
Non esiste il caso: \emph{tutti} i metodi sono virtuali. \\
\hline
Stesso nome, signature diversa nella derivata &
I metodi omonimi della base diventano \textbf{invisibili} (nascosti). &
Semplicemente \textbf{non} avviene override; i metodi restano visibili
(overloading). \\
\hline
Tipo di ritorno diverso, stessa signature &
Errore di compilazione (C++ ammette solo piccole variazioni:
\emph{covariant return types}). &
Errore, salvo tipo di ritorno covariante. \\
\hline
Impedire l'override & Non c'è \code{final} sui metodi (storicamente);
i metodi non-\code{virtual} semplicemente non si sovrascrivono. &
Parola chiave \code{final} (sul metodo o sull'intera classe). \\
\hline
\end{tabular}
\end{center}

\section{Sintassi nei quattro linguaggi}

\subsection{C++ --- virtuale esplicito}

In C++ un metodo è virtuale \textbf{solo se dichiarato tale}. I metodi non
virtuali usano early binding; se provi a "sovrascriverli" nella derivata li
stai solo \emph{nascondendo}.

\begin{lstlisting}[language=C++]
class B {
public:
    virtual void f() { std::cout << "B::f\n"; }  // virtuale
    void g()         { std::cout << "B::g\n"; }  // NON virtuale
};

class D : public B {
public:
    void f() override { std::cout << "D::f\n"; } // override (virtuale)
    void g()          { std::cout << "D::g\n"; } // HIDING, non override
};

int main() {
    B* p = new D();
    p->f();   // D::f  -> late binding (virtuale)
    p->g();   // B::g  -> early binding (tipo statico B)
    delete p;
}
\end{lstlisting}

\nt{
    In C++ (dal C++11) conviene sempre usare \code{override}: non cambia la
    semantica ma fa segnalare al compilatore se per errore \emph{non} stai
    davvero sovrascrivendo (es. signature sbagliata). \code{final} su un metodo
    o classe impedisce ulteriori override/derivazioni (introdotto in C++11,
    quindi \emph{oggi} un \code{final} esiste anche in C++, a differenza di
    quanto dicono le slide più vecchie).
}

\subsection{Java --- virtuale di default}

In Java \textbf{ogni} metodo non statico è virtuale. Non serve alcuna parola
chiave; l'override è sempre possibile (salvo \code{final}).

\begin{lstlisting}[language=Java]
class B {
    public void f() { System.out.println("B::f"); }
}
class D extends B {
    @Override
    public void f() { System.out.println("D::f"); }
}
public class Main {
    public static void main(String[] args) {
        B p = new D();
        p.f();   // D::f  -> sempre late binding
    }
}
\end{lstlisting}

\subsection{C --- niente OOP: puntatori a funzione}

In C non esistono classi né metodi virtuali. Il polimorfismo si \emph{emula}
mettendo puntatori a funzione dentro una \code{struct}: di fatto si costruisce
a mano quella \emph{tabella delle funzioni virtuali} che C++ e Java generano
automaticamente.

\begin{lstlisting}[language=C]
/* "vtable" costruita a mano */
typedef struct Animal {
    void (*voice)(struct Animal* self);   // puntatore a funzione = metodo virtuale
} Animal;

void dog_voice(Animal* self) { printf("bau\n"); }
void cat_voice(Animal* self) { printf("miao\n"); }

int main() {
    Animal a = { dog_voice };   // "istanza" con il suo comportamento
    Animal b = { cat_voice };
    a.voice(&a);   // bau   -> dispatch scelto dal puntatore a runtime
    b.voice(&b);   // miao
}
\end{lstlisting}

\nt{
    Questo è \emph{esattamente} ciò che il compilatore C++ fa sotto il cofano
    con la vtable (vedi sotto): la differenza è che in C++/Java è automatico e
    type-safe, in C è manuale ed è compito del programmatore mantenerlo
    corretto. Il prof cita proprio C e Pascal come linguaggi in cui il dispatch
    dinamico non esiste nativamente.
}

\subsection{Python --- tutto dinamico (duck typing)}

In Python il dispatch è \emph{sempre} dinamico e non serve dichiarare nulla: il
metodo viene cercato a runtime seguendo la \emph{Method Resolution Order} della
classe. Non conta la "classe del riferimento" (che non è nemmeno tipizzato).

\begin{lstlisting}[language=Python]
class Animal:
    def voice(self): print("???")

class Dog(Animal):
    def voice(self): print("bau")

class Cat(Animal):
    def voice(self): print("miao")

for x in (Dog(), Cat()):
    x.voice()      # bau, miao  -> sempre late binding
\end{lstlisting}

\dfn{Duck typing}{
    "If it walks like a duck and quacks like a duck, it's a duck". In Python
    non conta il \emph{tipo} dichiarato di un oggetto ma solo il fatto che
    \emph{possieda} il metodo/attributo richiesto. Il polimorfismo è quindi
    ancora più ampio di quello per ereditarietà: due classi \emph{non}
    imparentate sono interscambiabili se offrono gli stessi metodi.
}

\section{Come si implementa il late binding: la vtable}

\dfn{Virtual Function Table (VFT/vtable)}{
    Ogni classe che ha funzioni virtuali possiede una tabella (VFT) contenente
    i \textbf{puntatori al codice} di ciascuna sua funzione virtuale (VFP1,
    VFP2, \dots).
}

\dfn{Virtual Function Table Pointer (VFTP/vptr)}{
    Ogni \emph{oggetto} di una classe con funzioni virtuali contiene un
    puntatore nascosto (invisibile all'utente) alla vtable della sua classe.
}

\mprop{Passi del late binding a runtime}{
    Per eseguire la chiamata \code{p->VF2()} attraverso un puntatore \code{p}
    di tipo base:
    \begin{enumerate}
        \item dal puntatore si raggiunge l'oggetto;
        \item nell'oggetto si legge il VFTP (vptr);
        \item tramite il VFTP si raggiunge la vtable della classe \emph{reale};
        \item si cerca la voce giusta (per \code{VF2}, l'entry VFP2);
        \item si segue il puntatore trovato fino al codice e lo si esegue.
    \end{enumerate}
    Nota: in questo contesto "puntatore" significa puntatore C++ \emph{oppure}
    riferimento Java.
}

\nt{
    Java e Python usano meccanismi concettualmente identici: la JVM ha una
    tabella dei metodi per classe; CPython cerca il metodo nel \code{\_\_dict\_\_}
    della classe risalendo la MRO. Il costo dell'indirezione è lo stesso di cui
    parlano le slide.
}

\section{Virtual vs non-virtual: la questione dell'efficienza}

Se il polimorfismo è così comodo, perché non rendere tutti i metodi virtuali?
È la domanda con cui si spiega la differenza filosofica \textbf{Java vs C++}:

\begin{itemize}
    \item \textbf{Java} sceglie la semplicità: tutti i metodi virtuali, sempre.
    \item \textbf{C++} lascia la scelta all'utente (virtual/non-virtual). Il
    motivo è l'\textbf{efficienza}: il late binding costa (un'indirezione in
    più via vtable + impossibilità di fare \emph{inlining}).
\end{itemize}

\mprop{Quando l'overhead conta}{
    Il costo del late binding è di solito trascurabile. Diventa rilevante solo
    se un metodo \emph{semplice} viene chiamato \emph{moltissime} volte: in quei
    casi l'indirezione impedisce ottimizzazioni (inline) e può diventare un
    collo di bottiglia. Ambiti tipici: sistemi real-time, routine matematiche
    di base, programmazione di sistema di basso livello.
}

\begin{center}
\renewcommand{\arraystretch}{1.35}
\footnotesize
\begin{tabular}{|>{\raggedright\arraybackslash}p{3.4cm}|>{\raggedright\arraybackslash}p{3.4cm}|>{\raggedright\arraybackslash}p{3.4cm}|>{\raggedright\arraybackslash}p{3.0cm}|}
\hline
 & \textbf{C++} & \textbf{Java} & \textbf{Python} \\
\hline
Default & non virtuale & virtuale & virtuale \\
\hline
Costo di default & zero-overhead & indirezione vtable & lookup in dizionario \\
\hline
Ottimizzabile via inline & sì (non-virtual) & sì (JIT devirtualizza) & no (interprete) \\
\hline
Controllo al programmatore & massimo & minimo & nullo \\
\hline
\end{tabular}
\end{center}

\section{Esempio: tracciare l'output (dalle slide)}

Le slide propongono un esercizio classico d'esame: dato un insieme di classi
con metodi virtuali e non, capire quali chiamate compilano e quale output
producono. Il ragionamento chiave:

\mprop{Regola per l'esercizio}{
    \begin{itemize}
        \item Una chiamata \code{p->m()} attraverso un puntatore a \code{A}
        \textbf{compila} solo se \code{A} (o un suo antenato) dichiara un metodo
        con la signature di \code{m}: il compilatore guarda il \textbf{tipo
        statico} del puntatore per decidere se la chiamata è legale.
        \item Se il metodo è \code{virtual}, \emph{quale} versione venga
        eseguita dipende dal \textbf{tipo dinamico} (reale) dell'oggetto.
        \item Se il metodo \emph{non} è virtuale, si esegue la versione del
        \textbf{tipo statico}.
    \end{itemize}
}

\ex{Righe che non compilano}{
    Nell'esempio delle slide (classi \code{A..G}):
    \begin{itemize}
        \item \code{aG.f()} non compila se \code{G} non definisce un
        \code{f()} senza argomenti e \code{G} ha introdotto un \code{f(int)}
        che \emph{nasconde} gli altri \code{f}.
        \item \code{a->f(3)} con \code{a} di tipo \code{A*} non compila: il
        compilatore cerca \code{A::f(int)}, che non esiste, anche se \code{a}
        punta davvero a un oggetto derivato che ce l'ha (conta il tipo
        statico).
    \end{itemize}
    Sono errori dovuti al fatto che la \emph{legalità} della chiamata si decide
    sul tipo statico, mentre solo il \emph{dispatch} è dinamico.
}

%==============================================================
\chapter{Template e generics}
%==============================================================

Questo capitolo copre la lezione \emph{Templates}. Il problema di partenza è il
\textbf{riuso del codice}: molti programmi usano strutture dati generiche
(stack, code, liste). Duplicare il codice per ogni tipo di elemento (una coda
di clienti, una di messaggi, una di ordini) è insostenibile: ogni modifica va
replicata ovunque. La soluzione è \textbf{parametrizzare il codice rispetto al
tipo}.

\section{Il problema: container e tipizzazione}

\dfn{Genericità (type parameterization)}{
    Il meccanismo con cui una classe o una funzione viene \emph{parametrizzata
    rispetto a un tipo}, così da definire \emph{una sola volta} un algoritmo o
    una struttura dati che funziona per molti tipi diversi. Le due incarnazioni
    principali sono i \textbf{generics} (Java) e i \textbf{template} (C++).
}

Le slide legano la genericità al modello di tipizzazione:

\begin{itemize}
    \item Nei linguaggi a \textbf{tipizzazione dinamica} (Python) i container
    generici, sia \emph{omogenei} sia \emph{eterogenei}, si costruiscono
    banalmente: una lista può contenere qualunque cosa, non c'è un tipo da
    fissare.
    \item Nei linguaggi a \textbf{tipizzazione statica} (C, C++, Java) un
    container "cablato" a un tipo limita il riuso. Le soluzioni pre-generics
    erano:
    \begin{itemize}
        \item in \textbf{C}: \code{union}/variant record per l'eterogeneità (ma
        con un numero finito di alternative e cast "sporchi" in estrazione),
        oppure \code{void*} (che perde ogni type-safety);
        \item difficoltà a \emph{nascondere} l'implementazione: per iterare
        sul container serviva una variabile che rivelava la struttura interna.
    \end{itemize}
\end{itemize}

\nt{
    Ecco perché i template "nascono statici": servono proprio a recuperare, nei
    linguaggi a tipizzazione statica, quella flessibilità che i linguaggi
    dinamici hanno gratis --- ma \emph{senza} rinunciare al controllo dei tipi
    a compile-time.
}

\section{Generics e template: la differenza profonda}

\mprop{Template = generazione di codice a compile-time}{
    Un template è un \emph{modello} secondo cui viene generato codice reale.
    \begin{itemize}
        \item Quando il compilatore incontra la \emph{definizione} del template,
        \textbf{non genera codice}: fa i controlli possibili e rimanda il resto.
        \item Al punto di \textbf{istanziazione} (dove il template è usato con
        tipi concreti) il compilatore genera il codice \emph{specifico} per
        quei tipi. Il corpo eseguito non è generico: è del tutto specializzato.
        \item Il codice generato per una data combinazione di parametri è
        prodotto \textbf{una sola volta} e riusato.
    \end{itemize}
}

Questa è la distinzione capitale rispetto ai generics di Java:

\begin{center}
\renewcommand{\arraystretch}{1.4}
\footnotesize
\begin{tabular}{|>{\raggedright\arraybackslash}p{3.0cm}|>{\raggedright\arraybackslash}p{5.0cm}|>{\raggedright\arraybackslash}p{5.0cm}|}
\hline
 & \textbf{C++ template} & \textbf{Java generics} \\
\hline
Quando agiscono & compile-time: \textbf{monomorfizzazione} (una copia di codice per ogni tipo) & compile-time: \textbf{type erasure} (una sola copia, i parametri di tipo spariscono) \\
\hline
Codice prodotto & \code{vector<int>} e \code{vector<double>} sono classi \emph{diverse} e non imparentate & \code{List<Integer>} e \code{List<String>} condividono lo \emph{stesso} bytecode (\code{List}) \\
\hline
Tipi primitivi & funzionano direttamente & \textbf{vietati}: solo tipi riferimento (serve \code{Integer}, non \code{int}) \\
\hline
Type checking & incompleto sulla definizione, completo all'istanziazione & completo sulla definizione, poi cancellato \\
\hline
A runtime & tipo pienamente noto & tipo \textbf{cancellato} (\code{new T[]} impossibile, no \code{instanceof List<String>}) \\
\hline
Dimensione binario & può crescere (code bloat) & costante \\
\hline
Prestazioni & massime (nessun boxing, inline) & boxing/unboxing per i primitivi \\
\hline
\end{tabular}
\end{center}

\dfn{Monomorfizzazione}{
    Strategia (C++, Rust) in cui il compilatore genera una \emph{copia
    separata} del codice per ciascuna combinazione di parametri di tipo usata.
    Massime prestazioni, ma possibile aumento della dimensione del binario
    (\emph{code bloat}).
}

\dfn{Type erasure}{
    Strategia (Java) in cui i parametri di tipo servono solo al type-checker e
    vengono \emph{cancellati} dopo la compilazione: a runtime \code{List<String>}
    è semplicemente \code{List}. Vantaggio: un solo bytecode, compatibilità
    all'indietro. Svantaggio: informazioni di tipo non disponibili a runtime.
}

\nt{
    Le slide dicono che "i template non sono generics completi": non c'è
    typechecking completo sulla \emph{definizione}, perché finché i tipi sono
    ignoti il compilatore non può verificare che, per esempio, \code{r = 2 * t}
    sia lecito. In Java invece un generic è controllato \emph{subito} sulla
    definizione (grazie ai \emph{bounds}, es. \code{<T extends Comparable<T>>}).
}

\section{Function template}

\dfn{Function template}{
    Definisce una \emph{funzione generica} applicabile ad argomenti di tipi
    specificabili \emph{in seguito}, alla chiamata. Il controllo di consistenza
    dei tipi è rimandato alla prima chiamata (rilassamento del type checking).
}

\begin{lstlisting}[language=C++]
// dichiarazione: template + lista parametri tra <>
template <typename T1, typename T2>
void f(T1 t, T2& r) {
    r = 2 * t;   // legale solo per certi T1, T2: verificato all'istanziazione
}

template <class T>            // 'class' e 'typename' qui equivalenti
T maximum(T a, T b) {
    return a > b ? a : b;     // richiede che operator> sia definito su T
}

int main() {
    int x = 1, y = 2;
    f(y, x);                  // OK: deduce T1=int, T2=int&
    // f(anA, x);             // ERRORE: A non ha operator* con 2
    std::cout << maximum(10, 15);   // deduce T=int
}
\end{lstlisting}

\mprop{Deduzione dei parametri}{
    Nelle \emph{function} template il compilatore di solito \textbf{deduce} da
    solo i tipi guardando gli argomenti della chiamata. Se non ci riesce (o la
    chiamata è ambigua) bisogna specificarli a mano:
    \code{maximum<double>(a, b)}.
}

\subsection{Il classico tranello \code{char*}}

\begin{Example}{max su stringhe C}{}
    Con \code{maximum("Aladdin", "Jasmine")} il compilatore genera
    \code{maximum(char*, char*)}, che confronta gli \emph{indirizzi} dei due
    puntatori, non i contenuti! Risultato inatteso. Si risolve con una
    \textbf{specializzazione}:
\begin{lstlisting}[language=C++]
template <> char* maximum(char* a, char* b) {
    return strcmp(a, b) > 0 ? a : b;   // confronto vero delle stringhe
}
\end{lstlisting}
    In C++ moderno non si presenta con \code{std::string} (che sovraccarica
    \code{operator>}). È un ottimo esempio del prof su ``il template genera
    codice specifico che magari non fa ciò che credi''.
\end{Example}

\section{Class template}

\dfn{Class template}{
    Una classe parametrizzata rispetto a uno o più tipi. Sintassi analoga a
    quella dei function template.
}

\mprop{Istanziazione delle class template}{
    Differenza rispetto alle function template: per le \emph{classi} il
    compilatore \textbf{non deduce} i parametri (non ci sono argomenti da cui
    dedurli). Vanno \emph{sempre} specificati esplicitamente: \code{A<int>},
    \code{A<double>}.
}

\begin{lstlisting}[language=C++]
template <typename T>
class A {
    T x, y;
public:
    A() : x(0), y(0) {}
    T getSum() { return x + y; }
    void setXY(T a, T b) { x = a; y = b; }
};

int main() {
    A<int>    iA;    // genera la classe A<int>
    A<double> dA;    // genera la classe A<double>  (tipo DIVERSO da A<int>)
}
\end{lstlisting}

\mprop{Specializzazioni = classi indipendenti}{
    Ogni istanziazione con parametri diversi crea una nuova
    \textbf{specializzazione}. Due specializzazioni dello stesso template sono
    classi regolari e \textbf{completamente scorrelate}: \code{A<int>},
    \code{A<double>}, \code{A<B>} non hanno alcuna relazione. Non puoi assegnare
    un \code{A<int>} a un \code{A<double>} senza una conversione esplicita.
}

\subsection{Istanziazione implicita ed esplicita}

\begin{lstlisting}[language=C++]
Z<int> zi;                 // istanziazione IMPLICITA (all'uso)
template class Z<int>;     // istanziazione ESPLICITA (forzata dal programmatore)
\end{lstlisting}

\nt{
    Sottigliezza dalle slide: dichiarare solo un \emph{puntatore}
    \code{Z<int>* p;} \textbf{non} istanzia la classe (come dichiarare un
    puntatore a una struct incompleta). L'istanziazione scatta quando serve la
    definizione completa (es. creando un oggetto). Per questo, per i template,
    dichiarazione e definizione dei membri stanno di solito \emph{nello stesso
    header}: al punto di istanziazione il compilatore deve avere sotto gli occhi
    tutto il codice.
}

\section{Genericità nei quattro linguaggi}

\subsection{C --- nessun generics: macro, \code{void*}, \code{\_Generic}}

C non ha genericità nel type system. Tre approcci storici:

\begin{lstlisting}[language=C]
/* 1) Macro: sostituzione testuale del preprocessore (nessun type-check) */
#define MAX(a, b) ((a) > (b) ? (a) : (b))

/* 2) void*: container generico ma type-unsafe, richiede cast a mano */
typedef struct { void** data; int n; } List;

/* 3) qsort della libreria standard: generico via void* + comparatore */
int cmp_int(const void* a, const void* b) {
    return (*(const int*)a) - (*(const int*)b);
}
qsort(arr, n, sizeof(int), cmp_int);

/* 4) _Generic (C11): dispatch statico sul tipo, embrione di overloading */
#define abs_g(x) _Generic((x), int: abs, double: fabs)(x)
\end{lstlisting}

\nt{
    Questi meccanismi corrispondono alle "soluzioni pre-OO per i container"
    citate nelle slide: \code{void*} e \code{union} danno genericità ed
    eterogeneità, ma a costo di perdere il controllo dei tipi e di dover
    inserire "messy test and cast code" all'estrazione dei valori.
}

\subsection{C++ --- template (già visti)}

Monomorfizzazione a compile-time, massime prestazioni, base della
\emph{Standard Template Library} (STL: \code{vector}, \code{map}, \code{sort},
\dots).

\subsection{Java --- generics con erasure}

\begin{lstlisting}[language=Java]
class Cell<A> {                 // A = parametro di tipo
    private A value;
    Cell(A value) { this.value = value; }
    A get() { return value; }
}
// uso senza cast:
Cell<String> c = new Cell<>("First");
String s = c.get();             // nessun cast necessario

// bounded type: vincolo sui parametri (impossibile nei template classici)
static <T extends Comparable<T>> T max(T a, T b) {
    return a.compareTo(b) >= 0 ? a : b;
}
\end{lstlisting}

\nt{
    Conseguenze dell'erasure da ricordare: non puoi fare \code{new T()},
    \code{new T[n]}, né \code{o instanceof List<String>}; i generic \emph{non}
    accettano primitivi (\code{List<int>} è illegale, serve \code{List<Integer>}
    con boxing).
}

\subsection{Python --- duck typing e type hints}

Python non ha bisogno di generics a runtime: una funzione lavora su qualunque
oggetto che supporti le operazioni usate.

\begin{lstlisting}[language=Python]
def maximum(a, b):
    return a if a > b else b       # funziona per int, float, str, ...

# I "generics" esistono solo come ANNOTAZIONI statiche (typing),
# controllate da tool esterni (mypy), MAI a runtime:
from typing import TypeVar, Generic
T = TypeVar("T")
class Cell(Generic[T]):
    def __init__(self, value: T) -> None:
        self.value = value
    def get(self) -> T:
        return self.value
\end{lstlisting}

\mprop{Confronto sintetico sulla genericità}{
    \begin{itemize}
        \item \textbf{C}: genericità simulata (macro/\code{void*}), niente
        type-safety.
        \item \textbf{C++}: template $\to$ codice specializzato a compile-time,
        veloce ma può gonfiare il binario.
        \item \textbf{Java}: generics $\to$ un solo bytecode (erasure),
        type-safe a compile-time ma info perse a runtime.
        \item \textbf{Python}: non servono (duck typing); i type hint sono solo
        documentazione controllabile offline.
    \end{itemize}
}

%==============================================================
\chapter{Strategy pattern e costrutti abilitanti}
%==============================================================

Questo capitolo copre la lezione \emph{Focus on the strategy pattern in Java}.
Lo scopo del prof è mostrare come un \textbf{design pattern} si realizzi
concretamente nel codice, e come i costrutti del linguaggio (interfacce, classi
interne/anonime, polimorfismo, funzioni di prima classe) ne siano
l'infrastruttura.

\section{L'idea dello Strategy pattern}

\dfn{Strategy pattern}{
    Pattern comportamentale che \textbf{incapsula una famiglia di algoritmi}
    intercambiabili dietro una stessa interfaccia, permettendo di scegliere
    (anche a runtime) quale algoritmo usare, \emph{senza} modificare il codice
    cliente che lo invoca. La "strategia" è un oggetto (o una funzione) passato
    a chi la usa.
}

\mprop{Cosa abilita lo Strategy}{
    Lo Strategy si regge su uno di questi due meccanismi:
    \begin{enumerate}
        \item \textbf{polimorfismo per overriding + late binding}: la strategia
        è un oggetto che implementa un'interfaccia comune (Java, C++);
        \item \textbf{funzioni di prima classe}: la strategia è direttamente una
        funzione passata come argomento (C con puntatori a funzione, Python con
        funzioni/lambda, C++ con \code{std::function}).
    \end{enumerate}
    Il prof sottolinea che in Java lo Strategy è "polimorfismo via overriding e
    late binding", proprio \emph{perché} Java non ha puntatori a funzione.
}

\section{I costrutti Java che lo rendono elegante}

Le slide introducono una serie di costrutti Java (classi annidate, interne,
locali, anonime; interfacce con default/static) che servono a scrivere lo
Strategy in modo conciso.

\subsection{Classi annidate, interne, locali, anonime}

\dfn{Classi in scope interni}{
    Sia Java sia C++ permettono di definire una classe dentro un'altra. In C++
    è solo una questione di \emph{visibilità e naming}. In Java la distinzione è
    più profonda:
    \begin{itemize}
        \item \textbf{nested (static)}: dentro un'altra classe solo per
        organizzazione/naming; può usare solo i membri \emph{static}
        dell'esterna. Si dichiara con \code{static}.
        \item \textbf{inner (non static)}: ogni istanza è legata a
        un'\emph{istanza} della classe che la contiene e può accedere ai membri
        d'istanza di quest'ultima. Non può avere membri static.
        \item \textbf{local}: dichiarata dentro un \emph{blocco}/metodo; visibile
        solo lì; accede alle variabili locali \code{final}/effectively final.
        \item \textbf{anonymous}: classe locale \emph{senza nome}, dichiarata e
        istanziata in un'unica espressione; una sola istanza. Usata quando serve
        una classe una volta sola.
    \end{itemize}
}

\begin{lstlisting}[language=Java]
public class X {
    int instanceVar;
    static int staticVar;

    private class Inner {                 // INNER: legata all'istanza di X
        int y;
        Inner() { y = instanceVar + staticVar; }   // accede a entrambi
    }
    private static class Nested {         // NESTED: solo membri static
        int x;
        Nested() { x = staticVar; /* instanceVar -> ERRORE */ }
    }
}
\end{lstlisting}

In C++ esiste solo la versione "nested" (naming/visibilità); non c'è il legame
automatico con l'istanza dell'esterna:

\begin{lstlisting}[language=C++]
class X {
public:
    int instanceVar;
    static int staticVar;
    class Nested {
        int x;
        Nested() : x(staticVar) { /* instanceVar++ -> ERRORE */ }
    };
};
\end{lstlisting}

\subsubsection{Classi annidate in Python}

Python permette di definire una classe sia dentro un'altra classe sia dentro una
funzione, ma la \emph{semantica} è diversa da Java. La regola da ricordare:

\mprop{La classe annidata in Python $\approx$ nested \emph{static} di Java}{
    Una classe definita dentro un'altra è semplicemente un \textbf{attributo}
    della classe esterna (\code{Outer.Inner}). \textbf{Non} riceve alcun
    riferimento automatico all'istanza esterna: si comporta quindi come una
    \emph{static nested class} di Java, mai come una \emph{inner class}. La
    parola chiave \code{static} non esiste (e non serve) per le classi.
}

\nt{
    \textbf{Gotcha fondamentale (scope di classe non annidati).} Il corpo di una
    classe annidata \emph{non vede} i nomi definiti nel corpo della classe
    esterna: in Python solo gli scope di \emph{funzione} (via closure), globale e
    built-in fanno catena (LEGB), non quelli di classe. Verificato:
    \code{y = x} dentro \code{Inner} solleva \code{NameError}.
}

\begin{lstlisting}[language=Python]
class Outer:
    x = 1
    class Inner:
        # y = x        # -> NameError: lo scope della classe esterna NON e' visibile
        pass

o = Outer.Inner()      # Inner e' solo un attributo di Outer
# o non ha alcun riferimento automatico a un'istanza di Outer

# "inner class" con accesso all'esterna: si passa self ESPLICITAMENTE
class Container:
    def __init__(self): self.value = 42
    class Handle:
        def __init__(self, outer): self.outer = outer
        def show(self): return self.outer.value
    def make(self): return Container.Handle(self)   # legame esplicito
\end{lstlisting}

\paragraph{Classe locale (dentro una funzione).} A differenza degli scope di
classe, gli scope di \emph{funzione} in Python si annidano: una classe definita
dentro una funzione può leggere le variabili locali della funzione, sia nel
corpo della classe sia nei metodi (\textbf{closure}), senza bisogno di
\code{final} come in Java.

\begin{lstlisting}[language=Python]
def make_adder(n):
    class Adder:            # classe LOCALE
        base = n            # il corpo classe legge il locale n della funzione
        def add(self, x): return x + n   # metodo: closure su n
    return Adder()

make_adder(10).add(5)       # 15
\end{lstlisting}

\paragraph{Classi anonime: non esistono (e non servono).} Python non ha una
sintassi per le classi anonime. Al loro posto:
\begin{itemize}
    \item le \code{lambda} danno \emph{funzioni} anonime (una sola espressione);
    \item \code{type("Nome", (Base,), \{...\})} crea una classe al volo (ma vuole
    comunque un nome);
    \item nella stragrande maggioranza dei casi non serve nulla di tutto ciò:
    dove Java usa una classe anonima che implementa un'interfaccia a metodo
    singolo (es. \code{ActionListener}), Python passa direttamente una funzione
    (funzioni first-class + duck typing).
\end{itemize}

\nt{
    \textbf{Perché in Python le nested class sono rare.} Il meccanismo idiomatico
    di \emph{namespacing} in Python è il \textbf{modulo} (un file \code{.py}), non
    la classe annidata. Le classi annidate si usano quasi solo per legare
    strettamente un tipo di supporto alla sua classe (es. \code{Model.Meta} in
    Django) --- proprio l'uso "static nested" di Java.
}

\subsubsection{Classi annidate in C}

In C non esistono classi: l'analogo è la definizione di una \code{struct} dentro
un'altra \code{struct} (o dentro una funzione). Ma le regole di scope sono molto
diverse da C++/Java.

\mprop{In C il tag della struct annidata "leaka" nello scope esterno}{
    Contrariamente a C++ (dove \code{Inner} è \code{Outer::Inner}, scoped nella
    classe), in C il \emph{tag} di una struct annidata viene "issato" nello scope
    del file: diventa un tipo globale visibile ovunque. Non c'è quindi alcun
    vero \emph{namespacing}. (Verificato: il codice sotto compila in C, mentre in
    C++ \code{struct Inner leaked;} dà "incomplete type".)
}

\begin{lstlisting}[language=C]
struct Outer {
    int x;
    struct Inner { int y; } inner;   /* struct annidata */
};

struct Inner leaked;    /* C: OK, il tag Inner e' visibile a livello di file */
                        /* C++: ERRORE, servirebbe Outer::Inner */
\end{lstlisting}

\begin{itemize}
    \item \textbf{Struct locale}: una \code{struct} definita dentro una funzione
    è invece a scope di blocco (visibile solo lì), l'unico caso in cui C offre un
    vero incapsulamento del tipo.
    \item \textbf{Nessuna inner class}: per dare a una struct "annidata" accesso
    ai dati di quella esterna si passa un \emph{puntatore} all'esterna
    esplicitamente (come per lo Strategy e le funzioni virtuali in C).
    \item \textbf{Nessuna classe anonima}: attenzione, il C11 ha le
    \emph{anonymous struct/union}, ma sono \emph{membri} aggregati senza nome
    (servono per il layout, es. accedere a \code{s.a} invece di \code{s.u.a}), non
    hanno nulla a che vedere con le classi anonime di Java.
\end{itemize}

\subsubsection{Quadro comparativo}

\begin{center}
\renewcommand{\arraystretch}{1.35}
\footnotesize
\begin{tabular}{|>{\raggedright\arraybackslash}p{2.5cm}|>{\raggedright\arraybackslash}p{2.4cm}|>{\raggedright\arraybackslash}p{2.6cm}|>{\raggedright\arraybackslash}p{3.0cm}|>{\raggedright\arraybackslash}p{2.2cm}|}
\hline
\textbf{Concetto (Java)} & \textbf{Java} & \textbf{C++} & \textbf{Python} & \textbf{C} \\
\hline
Nested (static) & \code{static class} & \code{class} annidata (scoped) & classe annidata (default) & struct annidata (tag \emph{leaka}) \\
\hline
Inner (legata all'istanza) & \code{class} (ref. esterna implicita) & non esiste (ref. esplicito) & non esiste (passi \code{self}) & non esiste (passi puntatore) \\
\hline
Local (in funzione) & sì (inner, locali \code{final}) & sì & sì (closure sui locali) & sì (struct block-scoped) \\
\hline
Anonima & sì (\code{new I()\{...\}}) & lambda (C++11) & no (lambda / \code{type()}) & no (solo anon. struct \emph{member}) \\
\hline
Namespacing preferito & package + nested & namespace + nested & \textbf{modulo} & file + prefissi nomi \\
\hline
\end{tabular}
\end{center}

\subsection{Interfacce}

\dfn{Interfaccia (Java)}{
    Un tipo riferimento che dichiara \emph{signature} di metodi (astratti),
    costanti, e --- da Java 8 --- metodi \code{default} e \code{static}. Non ha
    variabili d'istanza e non può essere istanziata: le classi la
    \code{implements} fornendo il corpo dei metodi astratti.
}

\mprop{Perché le interfacce (e non l'ereditarietà multipla)}{
    Le interfacce permettono a classi \emph{non} imparentate di condividere un
    contratto comune \textbf{senza} una superclasse astratta condivisa. Java
    ammette \code{implements} di \emph{più} interfacce (ereditarietà multipla di
    \emph{tipo}) evitando le ambiguità dell'ereditarietà multipla di
    \emph{stato} tipica del C++ (il "diamond problem").
}

\dfn{Default method}{
    Metodo di interfaccia con un corpo (\code{default}). Serve a \textbf{far
    evolvere un'interfaccia} senza rompere le classi che già la implementano:
    aggiungere un metodo astratto le romperebbe, un \code{default} no (fornisce
    un'implementazione di fallback ereditabile).
}

\begin{lstlisting}[language=Java]
public interface DoIt {
    void doSomething(int i, double x);          // astratto
    int  doSomethingElse(String s);             // astratto
    default boolean didItWork(int i, double x, String s) {  // aggiunto dopo
        return true;                            // NON rompe le vecchie classi
    }
}
\end{lstlisting}

\mprop{Risoluzione dei conflitti tra default method}{
    Se una classe eredita più default con la stessa signature, Java applica:
    \begin{enumerate}
        \item i metodi \emph{d'istanza} di una superclasse battono i default
        d'interfaccia ("class wins");
        \item i metodi già sovrascritti da altri candidati vengono ignorati;
        \item se resta un conflitto irrisolubile, \textbf{errore di
        compilazione}: la classe \emph{deve} fare override esplicito, potendo
        invocare \code{Interfaccia.super.metodo()}.
    \end{enumerate}
}

\subsection{Classi anonime e listener}

Il caso d'uso principe è la gestione degli eventi (GUI): un \code{JButton} ha
bisogno di un oggetto che implementi \code{ActionListener} (interfaccia con il
solo metodo \code{actionPerformed}). Le slide mostrano tre modi di fornire la
strategia di reazione al click:

\begin{lstlisting}[language=Java]
// 1) la classe contenitore implementa direttamente l'interfaccia
public class SimpleCalc implements ActionListener {
    public void actionPerformed(ActionEvent e) { /* ... */ }
}

// 2) inner class dedicata (accede ai campi dell'esterna, parametrizzabile)
private class OperatorAction implements ActionListener {
    private int operator;
    OperatorAction(int op) { this.operator = op; }
    public void actionPerformed(ActionEvent e) { /* usa operator */ }
}

// 3) classe anonima: la strategia definita sul posto, dove serve
equalsButton.addActionListener(new ActionListener() {
    @Override public void actionPerformed(ActionEvent e) { /* ... */ }
});
\end{lstlisting}

\nt{
    La classe anonima \emph{sembra} passare "un pezzo di codice" come parametro,
    ma in realtà crea una classe a tutti gli effetti, ne istanzia un oggetto e
    passa \emph{quell'oggetto}. È il ponte concettuale verso le \textbf{lambda}
    (Java 8), che rendono lo Strategy ancora più conciso quando l'interfaccia ha
    un solo metodo (\emph{functional interface}):
    \code{equalsButton.addActionListener(e -> \{ ... \});}
}

\section{Lo Strategy nei quattro linguaggi}

Qui sta il cuore del confronto richiesto dal prof: \emph{lo stesso pattern},
quattro realizzazioni molto diverse a seconda di cosa il linguaggio offre.

\subsection{Java --- interfaccia + (anonima/lambda)}

\begin{lstlisting}[language=Java]
interface Strategy { int execute(int a, int b); }

class Context {
    private Strategy s;
    Context(Strategy s) { this.s = s; }
    int run(int a, int b) { return s.execute(a, b); }   // late binding
}

Context add = new Context((a, b) -> a + b);   // lambda = strategia
Context mul = new Context((a, b) -> a * b);
\end{lstlisting}

\subsection{C++ --- classe astratta, \code{std::function} o template}

C++ offre più alternative: la classica (classe base astratta + virtual),
oppure funzioni di prima classe con \code{std::function}, oppure la
\emph{policy-based design} con i template (strategia scelta a compile-time,
zero-overhead).

\begin{lstlisting}[language=C++]
// (a) classica: interfaccia = classe astratta con metodo virtuale puro
struct Strategy { virtual int execute(int, int) = 0; virtual ~Strategy() {} };
struct Add : Strategy { int execute(int a, int b) override { return a + b; } };

// (b) moderna: la strategia e' una funzione (first-class)
#include <functional>
std::function<int(int,int)> s = [](int a, int b){ return a + b; };

// (c) policy-based: strategia come parametro di template (compile-time)
template <class Policy>
struct Context { int run(int a, int b) { return Policy::execute(a, b); } };
\end{lstlisting}

\dfn{Funzione virtuale pura / classe astratta (C++)}{
    Un metodo dichiarato \code{= 0} è \emph{puro}: la classe che lo contiene è
    \textbf{astratta} (non istanziabile) e funge da interfaccia. È l'equivalente
    C++ dell'\code{interface} di Java; combinando più classi astratte si ottiene
    l'ereditarietà multipla di tipo.
}

\subsection{C --- puntatori a funzione (lo Strategy "originale")}

In C la strategia è semplicemente un \textbf{puntatore a funzione} passato al
codice cliente: è la forma più antica ed essenziale del pattern.

\begin{lstlisting}[language=C]
typedef int (*Strategy)(int, int);       // tipo "strategia"

int run(Strategy s, int a, int b) { return s(a, b); }

int add(int a, int b) { return a + b; }
int mul(int a, int b) { return a * b; }

int main() {
    printf("%d\n", run(add, 3, 4));      // 7
    printf("%d\n", run(mul, 3, 4));      // 12
}
\end{lstlisting}

\nt{
    Il prof osserva che in Java lo Strategy passa \emph{per forza} da
    interfacce e overriding "proprio perché non usiamo puntatori a funzione". In
    C accade l'opposto: non essendoci oggetti né override, il puntatore a
    funzione è \emph{l'unico} strumento --- ed è sufficiente.
}

\subsection{Python --- funzioni di prima classe (quasi banale)}

In Python le funzioni sono oggetti di prima classe: lo Strategy si riduce a
passare una funzione (o qualunque \emph{callable}). Non servono né interfacce né
classi.

\begin{lstlisting}[language=Python]
def run(strategy, a, b):
    return strategy(a, b)

run(lambda a, b: a + b, 3, 4)    # 7
run(lambda a, b: a * b, 3, 4)    # 12
\end{lstlisting}

\mprop{Perché lo stesso pattern cambia forma}{
    \begin{itemize}
        \item Dove ci sono \textbf{funzioni di prima classe} (C, Python, C++
        moderno) lo Strategy è "passa una funzione": leggero, poche righe.
        \item Dove il linguaggio è \textbf{puramente a oggetti senza puntatori a
        funzione} (Java classico) serve incapsulare la strategia in un
        \emph{oggetto} che implementa un'interfaccia --- da cui l'importanza di
        classi anonime e lambda per non appesantire il codice.
    \end{itemize}
    Il pattern è lo stesso; cambiano gli \emph{ingredienti} che il linguaggio
    mette a disposizione.
}

\section{Ereditarietà e polimorfismo: quadro comparativo}

Lo Strategy poggia sul polimorfismo, quindi conviene chiudere sistematizzando
come i quattro linguaggi gestiscono ereditarietà e interfacce.

\begin{center}
\renewcommand{\arraystretch}{1.35}
\footnotesize
\begin{tabular}{|>{\raggedright\arraybackslash}p{2.6cm}|>{\raggedright\arraybackslash}p{2.3cm}|>{\raggedright\arraybackslash}p{3.0cm}|>{\raggedright\arraybackslash}p{2.6cm}|>{\raggedright\arraybackslash}p{2.4cm}|}
\hline
\textbf{Aspetto} & \textbf{C} & \textbf{C++} & \textbf{Java} & \textbf{Python} \\
\hline
Ereditarietà & assente & multipla (di stato) & singola (di stato) & multipla (di stato) \\
\hline
Interfacce & assenti & classi astratte (virtual puro) & \code{interface} (multiple) & ABC / duck typing \\
\hline
Diamond problem & --- & sì $\to$ risolto con \code{virtual} inheritance & evitato: stato singolo, tipo multiplo & risolto dalla MRO (linearizzazione C3) \\
\hline
Metodo astratto & --- & \code{virtual f() = 0;} & \code{abstract} / metodo d'interfaccia & \code{@abstractmethod} \\
\hline
Strategia idiomatica & puntatore a funzione & virtual / \code{std::function} / template & interfaccia + lambda & funzione / callable \\
\hline
\end{tabular}
\end{center}

\dfn{Diamond problem ed ereditarietà virtuale}{
    Se \code{D} eredita da \code{B} e \code{C}, entrambe derivate da \code{A},
    in C++ \code{D} contiene \emph{due} copie di \code{A} (ambiguità sui membri
    di \code{A}). Si risolve con l'\textbf{ereditarietà virtuale}
    (\code{class B : virtual public A}), che fa condividere un'unica \code{A}.
    Java evita il problema alla radice: erediti lo \emph{stato} da una sola
    classe, e più interfacce (che stato non hanno). Python usa la
    \textbf{MRO C3} per linearizzare l'ordine di ricerca dei metodi.
}

%==============================================================
\chapter{Tabelle riassuntive per l'esame}
%==============================================================

Riepilogo "a colpo d'occhio" dei confronti chiesti dal professore. Queste
tabelle raccolgono, per ciascun argomento delle slide, la posizione dei quattro
linguaggi.

\section{Modello di esecuzione e tipizzazione}

\begin{center}
\renewcommand{\arraystretch}{1.35}
\footnotesize
\begin{tabular}{|>{\raggedright\arraybackslash}p{3.0cm}|>{\raggedright\arraybackslash}p{2.4cm}|>{\raggedright\arraybackslash}p{2.7cm}|>{\raggedright\arraybackslash}p{2.4cm}|>{\raggedright\arraybackslash}p{2.4cm}|}
\hline
 & \textbf{C} & \textbf{C++} & \textbf{Java} & \textbf{Python} \\
\hline
Compilazione & nativo & nativo & bytecode + JVM & bytecode + VM \\
\hline
Tipizzazione & statica debole & statica forte & statica forte & dinamica forte \\
\hline
Type-check tipi & compile-time & compile-time & compile-time & run-time \\
\hline
Memoria & manuale & manuale + RAII & GC & GC (refcount) \\
\hline
Velocità tipica & massima & massima & alta (JIT) & bassa \\
\hline
\end{tabular}
\end{center}

\section{Polimorfismo e dispatch (funzioni virtuali)}

\begin{center}
\renewcommand{\arraystretch}{1.35}
\footnotesize
\begin{tabular}{|>{\raggedright\arraybackslash}p{3.2cm}|>{\raggedright\arraybackslash}p{2.3cm}|>{\raggedright\arraybackslash}p{2.6cm}|>{\raggedright\arraybackslash}p{2.4cm}|>{\raggedright\arraybackslash}p{2.4cm}|}
\hline
 & \textbf{C} & \textbf{C++} & \textbf{Java} & \textbf{Python} \\
\hline
Dispatch di default & early (statico) & early (statico) & late (dinamico) & late (dinamico) \\
\hline
Metodo virtuale & simulato (ptr a funz.) & \code{virtual} esplicito & sempre & sempre \\
\hline
Bloccare l'override & --- & non-\code{virtual} / \code{final} (C++11) & \code{final} & convenzione / \code{\_\_} \\
\hline
Meccanismo & tabella manuale & vtable + vptr & tabella metodi JVM & lookup MRO / dict \\
\hline
Overhead & nullo & solo se virtual & sempre presente & interpretazione \\
\hline
\end{tabular}
\end{center}

\section{Genericità (template / generics)}

\begin{center}
\renewcommand{\arraystretch}{1.35}
\footnotesize
\begin{tabular}{|>{\raggedright\arraybackslash}p{3.2cm}|>{\raggedright\arraybackslash}p{2.3cm}|>{\raggedright\arraybackslash}p{2.7cm}|>{\raggedright\arraybackslash}p{2.5cm}|>{\raggedright\arraybackslash}p{2.4cm}|}
\hline
 & \textbf{C} & \textbf{C++} & \textbf{Java} & \textbf{Python} \\
\hline
Meccanismo & macro / \code{void*} & template & generics & duck typing \\
\hline
Quando & preproc. / runtime & compile-time & compile-time & runtime \\
\hline
Strategia & sostituzione testo & monomorfizzazione & type erasure & nessuna \\
\hline
Type-safe & no & sì (all'uso) & sì (poi erasa) & no (solo hint) \\
\hline
Primitivi & sì & sì & no (boxing) & n/a \\
\hline
Info a runtime & --- & sì & no & sì (dinamico) \\
\hline
\end{tabular}
\end{center}

\section{OOP: ereditarietà e interfacce (strategy)}

\begin{center}
\renewcommand{\arraystretch}{1.35}
\footnotesize
\begin{tabular}{|>{\raggedright\arraybackslash}p{3.2cm}|>{\raggedright\arraybackslash}p{2.2cm}|>{\raggedright\arraybackslash}p{2.7cm}|>{\raggedright\arraybackslash}p{2.6cm}|>{\raggedright\arraybackslash}p{2.4cm}|}
\hline
 & \textbf{C} & \textbf{C++} & \textbf{Java} & \textbf{Python} \\
\hline
OOP nativo & no & sì & sì & sì \\
\hline
Ereditarietà multipla & --- & sì (stato) & solo interfacce & sì (MRO) \\
\hline
Interfaccia & --- & classe astratta & \code{interface} & ABC / protocol \\
\hline
Strategy idiomatico & ptr a funzione & virtual / functional / template & interfaccia + lambda & callable \\
\hline
Classi anonime & --- & lambda (C++11) & sì & lambda \\
\hline
\end{tabular}
\end{center}

\section{Le domande "trappola" tipiche del prof}

\mprop{Cose da saper rispondere a voce}{
    \begin{itemize}
        \item \emph{Perché C++ ha \code{virtual} e Java no?} $\Rightarrow$
        efficienza / zero-overhead: C++ non fa pagare il late binding a chi non
        lo usa; Java preferisce l'uniformità.
        \item \emph{Differenza tra template C++ e generics Java?} $\Rightarrow$
        monomorfizzazione (codice per ogni tipo, runtime type noto) vs type
        erasure (un solo bytecode, tipo cancellato).
        \item \emph{Come si fa polimorfismo/strategy in C?} $\Rightarrow$
        puntatori a funzione (dentro struct = vtable manuale).
        \item \emph{Perché Python non ha bisogno di generics?} $\Rightarrow$
        duck typing: i tipi si controllano a runtime, una funzione lavora su
        qualunque oggetto "compatibile".
        \item \emph{Overriding vs overloading vs hiding?} $\Rightarrow$ stessa
        signature+virtuale (runtime) / stesso nome, signature diversa
        (compile-time) / (solo C++) omonimo nella derivata che nasconde la base.
        \item \emph{A cosa serve la vtable?} $\Rightarrow$ a realizzare il late
        binding: ogni oggetto ha un vptr alla tabella dei puntatori ai metodi
        virtuali della sua classe reale.
        \item \emph{Perché i default method nelle interfacce Java?}
        $\Rightarrow$ per far evolvere un'interfaccia senza rompere le classi
        che già la implementano.
    \end{itemize}
}

\vfill
\begin{center}
\itshape\small Fine degli appunti trasversali sui linguaggi.\\
Argomenti coperti: funzioni virtuali \& late binding, template \& generics,
strategy pattern --- confrontati su C, C++, Java, Python.
\end{center}

\end{document}
